\newMonth{Martuis}
\setrunningtitles{Martius}{March}

% ---- martyrology/mart03/mart0301.htm
\needspace{10\baselineskip}
\begin{paracol}{2}
\selectlanguage{latin}
\begin{center}{\color{gregoriocolor} Kaléndis Mártii. 
 Luna\dots\ }\end{center}
\switchcolumn
\selectlanguage{english}
\begin{center}{\color{gregoriocolor} The 1\textsuperscript{st} Day of March. 
 The\dots\ Day of the Moon.}\end{center}
\end{paracol}

\noindent\begin{tabularx}{\linewidth}{*{19}{>{\centering\arraybackslash}X}}
 \textcolor{gregoriocolor}{a} & \textcolor{gregoriocolor}{b} & \textcolor{gregoriocolor}{c} & \textcolor{gregoriocolor}{d} & \textcolor{gregoriocolor}{e} & \textcolor{gregoriocolor}{f} & \textcolor{gregoriocolor}{g} & \textcolor{gregoriocolor}{h} & \textcolor{gregoriocolor}{i} & \textcolor{gregoriocolor}{k} & \textcolor{gregoriocolor}{l} & \textcolor{gregoriocolor}{m} & \textcolor{gregoriocolor}{n} & \textcolor{gregoriocolor}{p} & \textcolor{gregoriocolor}{q} & \textcolor{gregoriocolor}{r} & \textcolor{gregoriocolor}{s} & \textcolor{gregoriocolor}{t} & \textcolor{gregoriocolor}{u} \\
 2 & 3 & 4 & 5 & 6 & 7 & 8 & 9 & 10 & 11 & 12 & 13 & 14 & 15 & 16 & 17 & 18 & 19 & 20 \\
\end{tabularx}
\vspace{0.5\baselineskip}
\noindent\begin{tabularx}{\linewidth}{*{12}{>{\centering\arraybackslash}X}}
 \textcolor{gregoriocolor}{A} & \textcolor{gregoriocolor}{B} & \textcolor{gregoriocolor}{C} & \textcolor{gregoriocolor}{D} & \textcolor{gregoriocolor}{E} & F & \textcolor{gregoriocolor}{F} & \textcolor{gregoriocolor}{G} & \textcolor{gregoriocolor}{H} & \textcolor{gregoriocolor}{M} & \textcolor{gregoriocolor}{N} & \textcolor{gregoriocolor}{P} \\
 21 & 22 & 23 & 24 & 25 & 26 & 25 & 26 & 27 & 28 & 29 & 1 \\
\end{tabularx}

\begin{paracol}{2}
\selectlanguage{latin}
\lettrine[lines=2]{R}{omæ} sanctórum Mártyrum ducentórum sexagínta, 
 quos jussit primo Cláudius, pro Christi nómine damnátos, extra portam 
 Saláriam arénam fódere, deínde in amphitheátro sagíttis mílitum intérfici.
\switchcolumn
\selectlanguage{english}
\lettrine[lines=2]{A}{t} Rome, two hundred and sixty holy martyrs condemned for the name of 
 Christ. Claudius ordered them to dig sand beyond the Salarian Gate, 
 then to have soldiers in the amphitheatre shoot them with arrows.
\switchcolumn*
\selectlanguage{latin}
Item natális sanctórum Mártyrum Leónis, Donáti, Abundántii, Nicéphori 
 et aliórum novem.
\switchcolumn
\selectlanguage{english}
Also, the birthday of the holy martyrs Leo, Donatus, Abundantius, Nicephorus, 
 and nine others.
\switchcolumn*
\selectlanguage{latin}
Massíliæ, in Gállia, sanctórum Mártyrum 
 Hermétis et Hadriáni.
\switchcolumn
\selectlanguage{english}
At Marseilles in France, the holy martyrs Hermes and Adrian.
\switchcolumn*
\selectlanguage{latin}
Heliópoli, apud Líbanum, sanctæ Eudóciæ 
 Mártyris, quæ, in persecutióne Trajáni, a Theódoto Epíscopo baptizáta et ad 
 certámen muníta, ibídem, Vincéntii Præsidis jussu percússa gládio, martyrii 
 corónam accépit.
\switchcolumn
\selectlanguage{english}
At Heliopolis, St. Eudocia, martyr in the persecution of Trajan. She 
 was baptized by Bishop Theodotus, and being fortified for the combat, was 
 put to the sword at the command of Vincent the governor, and thus received 
 the crown of martyrdom.
\switchcolumn*
\selectlanguage{latin}
Eódem die sanctæ Antonínæ Mártyris, quæ in 
 persecutióne Diocletiáni, cum Gentílium deos irrisísset, ídeo, post vários 
 cruciátus, in vase quodam inclúsa, in palúdem urbis Ceæ demérsa est.
\switchcolumn
\selectlanguage{english}
On the same day, St. Antonina, martyr. For deriding the gods of the 
 heathen, in the persecution of Diocletian, she was, after various torments, 
 shut up in a cask and drowned in a marsh near the city of Cea.
\switchcolumn*
\selectlanguage{latin}
Romæ natális sancti Felícis Papæ Tértii, qui 
 sancti Gregórii Magni átavus fuit; qui étiam (ut ipse Gregórius refert), 
 sanctæ Tharsíllæ nepti appárens, illam ad cæléstia regna vocávit.
\switchcolumn
\selectlanguage{english}
At Rome, the birthday of Pope St. Felix III, ancestor of St. Gregory the 
 Great, who relates of him that he appeared to St. Tharsilla, his niece, and 
 called her to the kingdom of heaven.
\switchcolumn*
\selectlanguage{latin}
Apud civitátem Werdam sancti Suitbérti Epíscopi, qui, témpore sancti Sérgii 
 Papæ Primi, apud Frísones, Bátavos et álios 
 Germániæ pópulos Evangélium prædicávit.
\switchcolumn
\selectlanguage{english}
At Kaiserswerdt, Bishop St. Swidbert, who, in the time of Pope Sergius, 
 preached the Gospel among the Frisians, Batavians, and other Germanic 
 peoples.
\switchcolumn*
\selectlanguage{latin}
Andégavi, in Gállia, sancti Albíni, Epíscopi et 
 Confessóris, viri præclaríssimæ virtútis et sanctitátis.
\switchcolumn
\selectlanguage{english}
At Angers in France, St. Albinus, bishop and confessor, a man of most 
 eminent virtue and piety.
\switchcolumn*
\selectlanguage{latin}
Apud Cenómanos, in Gállia, sancti Siviárdi 
 Abbátis.
\switchcolumn
\selectlanguage{english}
At Le Mans in France, St. Siviard, abbot.
\switchcolumn*
\selectlanguage{latin}
Perúsiæ Translátio sancti Herculáni, Epíscopi 
 et Mártyris, qui jussu Tótilæ, Gothórum Regis, decollátus est. Ipsíus 
 autem corpus ita cápiti unítum atque sanum, quadragésimo post abscissiónem 
 die (ut scribit sanctus Gregórius Papa), repértum est, ac si nulla ferri 
 incísio illud tetigísset.
\switchcolumn
\selectlanguage{english}
At Perugia, the transferral of the body of St. Herculanus, bishop and 
 martyr, who was beheaded by order of Totila, king of the Goths. Forty 
 days after the decapitation, Pope St. Gregory relates that the head had been 
 rejoined to the body as if it had never been touched by the sword.
\switchcolumn*
\selectlanguage{latin}
\end{paracol}

% \continueMonth{Martius}

% ---- martyrology/mart03/mart0302.htm
\needspace{10\baselineskip}
\begin{paracol}{2}
\selectlanguage{latin}
\begin{center}{\color{gregoriocolor} Sexto Nonas Mártii. 
 Luna\dots\ }\end{center}
\switchcolumn
\selectlanguage{english}
\begin{center}{\color{gregoriocolor} The 2\textsuperscript{nd} Day of March. 
 The\dots\ Day of the Moon.}\end{center}
\end{paracol}

\noindent\begin{tabularx}{\linewidth}{*{19}{>{\centering\arraybackslash}X}}
 \textcolor{gregoriocolor}{a} & \textcolor{gregoriocolor}{b} & \textcolor{gregoriocolor}{c} & \textcolor{gregoriocolor}{d} & \textcolor{gregoriocolor}{e} & \textcolor{gregoriocolor}{f} & \textcolor{gregoriocolor}{g} & \textcolor{gregoriocolor}{h} & \textcolor{gregoriocolor}{i} & \textcolor{gregoriocolor}{k} & \textcolor{gregoriocolor}{l} & \textcolor{gregoriocolor}{m} & \textcolor{gregoriocolor}{n} & \textcolor{gregoriocolor}{p} & \textcolor{gregoriocolor}{q} & \textcolor{gregoriocolor}{r} & \textcolor{gregoriocolor}{s} & \textcolor{gregoriocolor}{t} & \textcolor{gregoriocolor}{u} \\
 3 & 4 & 5 & 6 & 7 & 8 & 9 & 10 & 11 & 12 & 13 & 14 & 15 & 16 & 17 & 18 & 19 & 20 & 21 \\
\end{tabularx}
\vspace{0.5\baselineskip}
\noindent\begin{tabularx}{\linewidth}{*{12}{>{\centering\arraybackslash}X}}
 \textcolor{gregoriocolor}{A} & \textcolor{gregoriocolor}{B} & \textcolor{gregoriocolor}{C} & \textcolor{gregoriocolor}{D} & \textcolor{gregoriocolor}{E} & F & \textcolor{gregoriocolor}{F} & \textcolor{gregoriocolor}{G} & \textcolor{gregoriocolor}{H} & \textcolor{gregoriocolor}{M} & \textcolor{gregoriocolor}{N} & \textcolor{gregoriocolor}{P} \\
 22 & 23 & 24 & 25 & 26 & 27 & 26 & 27 & 28 & 29 & 1 & 2 \\
\end{tabularx}

\begin{paracol}{2}
\selectlanguage{latin}
\lettrine[lines=2]{R}{omæ,} via Latína, sanctórum Mártyrum Jovíni et 
 Basiléi, qui passi sunt sub Valeriáno et Galliéno Imperatóribus.
\switchcolumn
\selectlanguage{english}
\lettrine[lines=2]{A}{t} Rome, on the Latin Way, the holy martyrs Jovinus and Basileus, who 
 suffered under Emperors Valerian and Gallienus.
\switchcolumn*
\selectlanguage{latin}
Item Romæ plurimórum sanctórum Mártyrum, qui 
 sub Alexándro Imperatóre et Ulpiáno Præfécto, diu cruciáti, ad extrémum 
 capitáli senténtia damnáti sunt.
\switchcolumn
\selectlanguage{english}
Also at Rome, under Emperor Alexander and the prefect Ulpian, many holy 
 martyrs who were a long time tortured and condemned to capital punishment.
\switchcolumn*
\selectlanguage{latin}
Cæsaréæ, in Cappadócia, sanctórum Mártyrum 
 Lúcii Epíscopi, Absalónis et Lórgii.
\switchcolumn
\selectlanguage{english}
At Caesarea, in Cappadocia, the holy martyrs Lucius, bishop, Absalon, and 
 Lorgius.
\switchcolumn*
\selectlanguage{latin}
In Portu Románo sanctórum Mártyrum Pauli, Heráclíi, Secundíllæ 
 et Januáriæ.
\switchcolumn
\selectlanguage{english}
At Porto, near Rome, the holy martyrs Paul, Heraclius, Secundilla, and 
 Januaria.
\switchcolumn*
\selectlanguage{latin}
In Campánia commemorátio sanctórum octogínta Mártyrum, qui, cum nollent 
 carnes immolátas comédere nec caput capræ 
 adoráre, a Longobárdis sævíssime cæsi sunt.
\switchcolumn
\selectlanguage{english}
In Campania, the commemoration of eighty holy martyrs, who were barbarously 
 killed by the Lombards because they would not eat flesh that had been 
 offered to the idols, nor would they adore the head of a goat.
\switchcolumn*
\selectlanguage{latin}
Lichféldiæ, in Anglia, sancti Ceáddæ, Epíscopi 
 Merciórum et Lindisfarnórum, cujus præcláras virtútes sanctus Beda 
 Venerábilis commémorat.
\switchcolumn
\selectlanguage{english}
At Lichfield in England, St. Chad, bishop of Mercia and Lindisfarne, whose 
 excellent virtues are mentioned by St. Venerable Bede.
\switchcolumn*
\selectlanguage{latin}
\end{paracol}


% ---- martyrology/mart03/mart0303.htm
\needspace{10\baselineskip}
\begin{paracol}{2}
\selectlanguage{latin}
\begin{center}{\color{gregoriocolor} Quinto Nonas Mártii. 
 Luna\dots\ }\end{center}
\switchcolumn
\selectlanguage{english}
\begin{center}{\color{gregoriocolor} The 3\textsuperscript{rd} Day of March. 
 The\dots\ Day of the Moon.}\end{center}
\end{paracol}

\noindent\begin{tabularx}{\linewidth}{*{19}{>{\centering\arraybackslash}X}}
 \textcolor{gregoriocolor}{a} & \textcolor{gregoriocolor}{b} & \textcolor{gregoriocolor}{c} & \textcolor{gregoriocolor}{d} & \textcolor{gregoriocolor}{e} & \textcolor{gregoriocolor}{f} & \textcolor{gregoriocolor}{g} & \textcolor{gregoriocolor}{h} & \textcolor{gregoriocolor}{i} & \textcolor{gregoriocolor}{k} & \textcolor{gregoriocolor}{l} & \textcolor{gregoriocolor}{m} & \textcolor{gregoriocolor}{n} & \textcolor{gregoriocolor}{p} & \textcolor{gregoriocolor}{q} & \textcolor{gregoriocolor}{r} & \textcolor{gregoriocolor}{s} & \textcolor{gregoriocolor}{t} & \textcolor{gregoriocolor}{u} \\
 4 & 5 & 6 & 7 & 8 & 9 & 10 & 11 & 12 & 13 & 14 & 15 & 16 & 17 & 18 & 19 & 20 & 21 & 22 \\
\end{tabularx}
\vspace{0.5\baselineskip}
\noindent\begin{tabularx}{\linewidth}{*{12}{>{\centering\arraybackslash}X}}
 \textcolor{gregoriocolor}{A} & \textcolor{gregoriocolor}{B} & \textcolor{gregoriocolor}{C} & \textcolor{gregoriocolor}{D} & \textcolor{gregoriocolor}{E} & F & \textcolor{gregoriocolor}{F} & \textcolor{gregoriocolor}{G} & \textcolor{gregoriocolor}{H} & \textcolor{gregoriocolor}{M} & \textcolor{gregoriocolor}{N} & \textcolor{gregoriocolor}{P} \\
 23 & 24 & 25 & 26 & 27 & 28 & 27 & 28 & 29 & 1 & 2 & 3 \\
\end{tabularx}

\begin{paracol}{2}
\selectlanguage{latin}
\lettrine[lines=2]{C}{æsaréæ,} in Palæstína, sanctórum Mártyrum 
 Maríni mílitis, et Astérii Senatóris, in persecutióne Valeriáni. Horum 
 prior, cum accusátus esset a commilitiónibus ut Christiánus, et, 
 interrogátus a Júdice, se Christiánum esse voce claríssima testarétur, 
 martyrii corónam abscissióne cápitis accépit; cumque Astérius corpus 
 Mártyris, cápite truncátum, subjéctis húmeris et substráta veste, qua 
 induebátur, excíperet, honórem quem Martyrii détulit, contínuo et ipse 
 Martyr accépit.
\switchcolumn
\selectlanguage{english}
\lettrine[lines=2]{A}{t} Caesarea in Palestine, during the persecution of Valerian, the holy 
 martyrs Marinus, soldier, and Asterius, senator. The former was 
 examined by the judge on the charge laid against him by his fellow soldiers 
 of being a Christian, and as he admitted the accusation in a firm tone of 
 voice, he was beheaded, and thus received the crown of martyrdom. His 
 mutilated body was taken by Asterius on his own shoulders, and wrapped in 
 the garment which he himself wore. This service at once gained for 
 Asterius the palm of martyrdom as a reward for the honour which he had given 
 to a martyr.
\switchcolumn*
\selectlanguage{latin}
Calagúrri, in Hispánia, natális sanctórum Mártyrum Hemitérii et Cheledónii 
 fratrum, qui, cum apud Legiónem, Gallǽciæ 
 urbem, in castris militárent, ambo, exsurgénte persecutiónis procélla, pro 
 confessióne nóminis Christi, Calagúrrim usque profécti, ibi, plúribus 
 torméntis afflícti, martyrio coronáti sunt.
\switchcolumn
\selectlanguage{english}
At Calahorra in Spain, the birthday of the holy martyrs Hermiterius and 
 Cheledonius, soldiers in the army at Leon, a city of Galicia. Upon the 
 approach of persecution they went to Calahorra in order to confess the name 
 of Christ, and after enduring many torments there, they were crowned with 
 martyrdom.
\switchcolumn*
\selectlanguage{latin}
Eódem die pássio sanctórum Felícis, Lucióli, Fortunáti, Márciæ 
 et Sociórum.
\switchcolumn
\selectlanguage{english}
The same day, the passion of the Saints Felix, Luciolus, Fortunatus, Marcia, 
 and their companions.
\switchcolumn*
\selectlanguage{latin}
Item sanctórum mílitum Cleoníci, Eutrópii et 
 Basilísci, qui, in persecutióne Maximiáni, sub Asclepíade Præside, crucis 
 supplício felíciter triumphárunt.
\switchcolumn
\selectlanguage{english}
Also, the sainted soldiers Cleonicus, Eutropius, and Basiliscus, who 
 gloriously triumphed by death on the cross under the governor Asclepias 
 during the persecution of Maximian.
\switchcolumn*
\selectlanguage{latin}
Bríxiæ sancti Titiáni, Epíscopi et Confessóris.
\switchcolumn
\selectlanguage{english}
At Brescia, St. Titian, bishop and confessor.
\switchcolumn*
\selectlanguage{latin}
Bambérgæ sanctæ Cunegúndis Augústæ, quæ, sancto 
 Henríco Primo, Romanórum Imperatóri, nupta, perpétuam virginitátem, ipso 
 annuénte, servávit; ac, bonórum óperum méritis cumuláta, sancto fine quiévit, 
 et post óbitum miráculis cláruit.
\switchcolumn
\selectlanguage{english}
At Bamberg, Empress St. Cunegunda, who preserved her virginity with the 
 consent of her husband, Emperor Henry I. She completed a life rich in 
 meritorious good works with a holy death, and afterward worked many 
 miracles.
\switchcolumn*
\selectlanguage{latin}
\end{paracol}


% ---- martyrology/mart03/mart0304.htm
\needspace{10\baselineskip}
\begin{paracol}{2}
\selectlanguage{latin}
\begin{center}{\color{gregoriocolor} Quarto Nonas Mártii. 
 Luna\dots\ }\end{center}
\switchcolumn
\selectlanguage{english}
\begin{center}{\color{gregoriocolor} The 4\textsuperscript{th} Day of March. 
 The\dots\ Day of the Moon.}\end{center}
\end{paracol}

\noindent\begin{tabularx}{\linewidth}{*{19}{>{\centering\arraybackslash}X}}
 \textcolor{gregoriocolor}{a} & \textcolor{gregoriocolor}{b} & \textcolor{gregoriocolor}{c} & \textcolor{gregoriocolor}{d} & \textcolor{gregoriocolor}{e} & \textcolor{gregoriocolor}{f} & \textcolor{gregoriocolor}{g} & \textcolor{gregoriocolor}{h} & \textcolor{gregoriocolor}{i} & \textcolor{gregoriocolor}{k} & \textcolor{gregoriocolor}{l} & \textcolor{gregoriocolor}{m} & \textcolor{gregoriocolor}{n} & \textcolor{gregoriocolor}{p} & \textcolor{gregoriocolor}{q} & \textcolor{gregoriocolor}{r} & \textcolor{gregoriocolor}{s} & \textcolor{gregoriocolor}{t} & \textcolor{gregoriocolor}{u} \\
 5 & 6 & 7 & 8 & 9 & 10 & 11 & 12 & 13 & 14 & 15 & 16 & 17 & 18 & 19 & 20 & 21 & 22 & 23 \\
\end{tabularx}
\vspace{0.5\baselineskip}
\noindent\begin{tabularx}{\linewidth}{*{12}{>{\centering\arraybackslash}X}}
 \textcolor{gregoriocolor}{A} & \textcolor{gregoriocolor}{B} & \textcolor{gregoriocolor}{C} & \textcolor{gregoriocolor}{D} & \textcolor{gregoriocolor}{E} & F & \textcolor{gregoriocolor}{F} & \textcolor{gregoriocolor}{G} & \textcolor{gregoriocolor}{H} & \textcolor{gregoriocolor}{M} & \textcolor{gregoriocolor}{N} & \textcolor{gregoriocolor}{P} \\
 24 & 25 & 26 & 27 & 28 & 29 & 28 & 29 & 1 & 2 & 3 & 4 \\
\end{tabularx}

\begin{paracol}{2}
\selectlanguage{latin}
\lettrine[lines=2]{V}{ilnæ,} in Lithuánia, beáti Casimíri Confessóris, 
 e Casimíro Rege progéniti; quem Leo Décimus, Románus Póntifex, in Sanctórum 
 númerum rétulit.
\switchcolumn
\selectlanguage{english}
\lettrine[lines=2]{A}{t} Vilnius in Lithuania, blessed Casimir, confessor, the son of King Casimir, 
 whom Pope Leo X inscribed in the roll of the saints.
\switchcolumn*
\selectlanguage{latin}
Romæ, via Appia, natális sancti Lúcii Primi, 
 Papæ et Mártyris; qui, primo in persecutióne Valeriáni ob Christi fidem 
 exsílio relegátus, et póstmodum divíno nutu ad Ecclésiam suam redíre 
 permíssus, tandem, cum plúrimum advérsus Novatiános laborásset, cápitis 
 obtruncatióne martyrium complévit. Eum vero sanctus Cypriánus summis 
 láudibus celebrávit.
\switchcolumn
\selectlanguage{english}
At Rome, on the Appian Way, during the persecution of Valerian, the birthday 
 of St. Lucius, pope and martyr, who was first exiled for the faith of 
 Christ, but being permitted by divine Providence to return to his church, 
 after labouring long against the Novatians, he suffered martyrdom by 
 beheading. His praises have been published by St. Cyprian.
\switchcolumn*
\selectlanguage{latin}
Nicomedíæ sancti Hadriáni Mártyris, cum áliis 
 vigínti tribus, qui omnes, sub Diocletiáno Imperatóre, martyrium crurifrágio 
 consummárunt. Eórum relíquiæ, a Christiánis Byzántium delátæ, 
 reverénti honóre sepúltæ fuérunt; inde póstea sancti Hadriáni corpus Romam 
 translátum fuit sexto Idus Septémbris, quo die festívitas ejus potíssimum 
 celebrátur.
\switchcolumn
\selectlanguage{english}
At Nicomedia, in the reign of Emperor Diocletian, the martyr St. Adrian and 
 twenty-three others, who endured martyrdom by having their limbs crushed. 
 Their remains were taken to Byzantium by the Christians, and buried there 
 with reverence and honour. Afterwards the body of St. Adrian was 
 transferred to Rome on the 8th of September, on which day his feast is 
 observed.
\switchcolumn*
\selectlanguage{latin}
Romæ, via Appia, sanctórum Mártyrum nongentórum, 
 qui pósiti sunt in cœmetério ad sanctam Cæcíliam.
\switchcolumn
\selectlanguage{english}
At Rome, on the Appian Way, nine hundred holy martyrs who were buried in the 
 cemetery of St. Cecilia.
\switchcolumn*
\selectlanguage{latin}
Apud Chersonésum pássio sanctórum Episcopórum Basilíi, 
 Eugénii, Agathodóri, Elpídii, Æthérii, Capitónis, Ephræm, Néstoris et 
 Arcádii.
\switchcolumn
\selectlanguage{english}
In Chersonesus, the passion of the saintly bishops, Basil, Eugene, 
 Agathodorus, Elpidius, Aetherius, Capito, Ephrem, Nestor, and Arcadius.
\switchcolumn*
\selectlanguage{latin}
Eódem die sancti Caji Palatíni, in mare demérsi, et aliórum vigínti septem.
\switchcolumn
\selectlanguage{english}
On the same day, St. Caius Palatinus and twenty-seven others who were cast 
 into the sea.
\switchcolumn*
\selectlanguage{latin}
Item pássio sanctórum Archelái, Cyrílli et Phótii.
\switchcolumn
\selectlanguage{english}
Also, the passion of Saints Archelaus, Cyril and Photius.
\switchcolumn*
\selectlanguage{latin}
\end{paracol}


% ---- martyrology/mart03/mart0305.htm
\needspace{10\baselineskip}
\begin{paracol}{2}
\selectlanguage{latin}
\begin{center}{\color{gregoriocolor} Tértio Nonas Mártii. 
 Luna\dots\ }\end{center}
\switchcolumn
\selectlanguage{english}
\begin{center}{\color{gregoriocolor} The 5\textsuperscript{th} Day of March. 
 The\dots\ Day of the Moon.}\end{center}
\end{paracol}

\noindent\begin{tabularx}{\linewidth}{*{19}{>{\centering\arraybackslash}X}}
 \textcolor{gregoriocolor}{a} & \textcolor{gregoriocolor}{b} & \textcolor{gregoriocolor}{c} & \textcolor{gregoriocolor}{d} & \textcolor{gregoriocolor}{e} & \textcolor{gregoriocolor}{f} & \textcolor{gregoriocolor}{g} & \textcolor{gregoriocolor}{h} & \textcolor{gregoriocolor}{i} & \textcolor{gregoriocolor}{k} & \textcolor{gregoriocolor}{l} & \textcolor{gregoriocolor}{m} & \textcolor{gregoriocolor}{n} & \textcolor{gregoriocolor}{p} & \textcolor{gregoriocolor}{q} & \textcolor{gregoriocolor}{r} & \textcolor{gregoriocolor}{s} & \textcolor{gregoriocolor}{t} & \textcolor{gregoriocolor}{u} \\
 6 & 7 & 8 & 9 & 10 & 11 & 12 & 13 & 14 & 15 & 16 & 17 & 18 & 19 & 20 & 21 & 22 & 23 & 24 \\
\end{tabularx}
\vspace{0.5\baselineskip}
\noindent\begin{tabularx}{\linewidth}{*{12}{>{\centering\arraybackslash}X}}
 \textcolor{gregoriocolor}{A} & \textcolor{gregoriocolor}{B} & \textcolor{gregoriocolor}{C} & \textcolor{gregoriocolor}{D} & \textcolor{gregoriocolor}{E} & F & \textcolor{gregoriocolor}{F} & \textcolor{gregoriocolor}{G} & \textcolor{gregoriocolor}{H} & \textcolor{gregoriocolor}{M} & \textcolor{gregoriocolor}{N} & \textcolor{gregoriocolor}{P} \\
 25 & 26 & 27 & 28 & 29 & 30 & 29 & 1 & 2 & 3 & 4 & 5 \\
\end{tabularx}

\begin{paracol}{2}
\selectlanguage{latin}
\lettrine[lines=2]{A}{ntiochíæ} natális sancti Phocæ Mártyris, qui, 
 post multas, quas pro nómine Redemptoris passus est, injúrias, quáliter de 
 antíquo illo serpénte triumpháverit, hódie quoque pópulis eo miráculo 
 declarátur, quod, si quíspiam a serpénte morsus fúerit, hic, ut jánuam 
 Basílicæ Mártyris credens attígerit, conféstim, evacuáta venéni virtúte, 
 sanátur.
\switchcolumn
\selectlanguage{english}
\lettrine[lines=2]{A}{t} Antioch, the birthday of the martyr St. Phocas, who triumphed over the 
 ageless Serpent after many injuries which he suffered for the Name of the 
 Redeemer. That triumph is still manifested to the people in our day, 
 for if any one stung by a snake touches with faith the door of the martyr's 
 basilica, the power of the venom disappears, and he is immediately cured.
\switchcolumn*
\selectlanguage{latin}
Cæsaréæ, in Palæstína, sancti Hadriáni Mártyris, 
 qui in persecutióne Diocletiáni Imperatóris, jussu Firmiliáni Præsidis, 
 prius ob Christi fidem leóni objéctus, deínde gládio jugulátus, martyrii 
 corónam accépit.
\switchcolumn
\selectlanguage{english}
At Caesarea in Palestine, in the persecution of Diocletian, the martyr St. 
 Adrian. He was first exposed to a lion for the faith of Christ, and 
 then slain with the sword by order of the governor Firmilian, and thus 
 received the crown of martyrdom.
\switchcolumn*
\selectlanguage{latin}
Eódem die pássio sanctórum Eusébii Palatini, et aliórum novem Mártyrum.
\switchcolumn
\selectlanguage{english}
The same day, the passion of the holy martyrs Eusebius Palatinus and nine 
 others.
\switchcolumn*
\selectlanguage{latin}
Cæsearéæ, in Palæstína, sancti Theóphili 
 Epíscopi, qui sub Sevéro Príncipe, sapiéntia et vitæ integritáte conspícuus, 
 emícuit.
\switchcolumn
\selectlanguage{english}
At Caesarea in Palestine, in the time of Emperor Severus, St. Theophilus, 
 bishop, who was conspicuous for his wisdom and the purity of his life.
\switchcolumn*
\selectlanguage{latin}
Ad ripam Jordánis, item in Palæstína, sancti 
 Gerásimi, Anachorétæ et Abbátis; qui témpore Zenónis Imperatóris flóruit.
\switchcolumn
\selectlanguage{english}
Also in Palestine, on the banks of the Jordan, the anchoret St. Gerasimus, 
 who lived in the time of Emperor Zeno.
\switchcolumn*
\selectlanguage{latin}
Neápoli, in Campánia, deposítio sancti Joánnis-Joséphi a Cruce, Sacerdótis 
 ex Ordine Minórum et Confessóris, qui, sanctórum Ordini Seráphico insígne 
 decus áddidit, atque a Gregório Papa Décimo sexto in Sanctórum cánonem 
 est relátus.
\switchcolumn
\selectlanguage{english}
At Naples, in Campania, the death of St. John Joseph of the Cross, priest of 
 the Order of Friars Minor, and confessor. By emulating the virtues of 
 St. Francis of Assisi and of St. Peter Alcantara, he added great glory to 
 the Seraphic Order. He was canonized by Pope Gregory XVI.
\switchcolumn*
\selectlanguage{latin}
\end{paracol}


% ---- martyrology/mart03/mart0306.htm
\needspace{10\baselineskip}
\begin{paracol}{2}
\selectlanguage{latin}
\begin{center}{\color{gregoriocolor} Prídie Nonas Mártii. 
 Luna\dots\ }\end{center}
\switchcolumn
\selectlanguage{english}
\begin{center}{\color{gregoriocolor} The 6\textsuperscript{th} Day of March. 
 The\dots\ Day of the Moon.}\end{center}
\end{paracol}

\noindent\begin{tabularx}{\linewidth}{*{19}{>{\centering\arraybackslash}X}}
 \textcolor{gregoriocolor}{a} & \textcolor{gregoriocolor}{b} & \textcolor{gregoriocolor}{c} & \textcolor{gregoriocolor}{d} & \textcolor{gregoriocolor}{e} & \textcolor{gregoriocolor}{f} & \textcolor{gregoriocolor}{g} & \textcolor{gregoriocolor}{h} & \textcolor{gregoriocolor}{i} & \textcolor{gregoriocolor}{k} & \textcolor{gregoriocolor}{l} & \textcolor{gregoriocolor}{m} & \textcolor{gregoriocolor}{n} & \textcolor{gregoriocolor}{p} & \textcolor{gregoriocolor}{q} & \textcolor{gregoriocolor}{r} & \textcolor{gregoriocolor}{s} & \textcolor{gregoriocolor}{t} & \textcolor{gregoriocolor}{u} \\
 7 & 8 & 9 & 10 & 11 & 12 & 13 & 14 & 15 & 16 & 17 & 18 & 19 & 20 & 21 & 22 & 23 & 24 & 25 \\
\end{tabularx}
\vspace{0.5\baselineskip}
\noindent\begin{tabularx}{\linewidth}{*{12}{>{\centering\arraybackslash}X}}
 \textcolor{gregoriocolor}{A} & \textcolor{gregoriocolor}{B} & \textcolor{gregoriocolor}{C} & \textcolor{gregoriocolor}{D} & \textcolor{gregoriocolor}{E} & F & \textcolor{gregoriocolor}{F} & \textcolor{gregoriocolor}{G} & \textcolor{gregoriocolor}{H} & \textcolor{gregoriocolor}{M} & \textcolor{gregoriocolor}{N} & \textcolor{gregoriocolor}{P} \\
 26 & 27 & 28 & 29 & 30 & 1 & 1 & 2 & 3 & 4 & 5 & 6 \\
\end{tabularx}

\begin{paracol}{2}
\selectlanguage{latin}
\lettrine[lines=2]{S}{anctárum} Perpétuæ et Felicitátis Mártyrum, quæ 
 sequénti die gloriósam martyrii corónam a Dómino recepérunt.
\switchcolumn
\selectlanguage{english}
\lettrine[lines=2]{S}{aints} Perpetua and Felicity, who, on the day following this, received from 
 the Lord the glorious crown of martyrdom.
\switchcolumn*
\selectlanguage{latin}
Dertónæ sancti Marciáni, Epíscopi et Mártyris, 
 qui sub Trajáno, pro Christi glória occísus, coronátur.
\switchcolumn
\selectlanguage{english}
At Tortona, St. Marcian, bishop and martyr, who was put to death for the 
 sake of Christ by Trajan, and thereby received the crown of immortality.
\switchcolumn*
\selectlanguage{latin}
Nicomedíæ natális sanctórum Mártyrum Victóris 
 et Victoríni, qui, per triénnium, cum Claudiáno et uxóre ejus Bassa, 
 torméntis multis afflícti et retrúsi in cárcerem, ibídem vitæ suæ cursum 
 implevérunt.
\switchcolumn
\selectlanguage{english}
At Nicomedia, the birthday of the holy martyrs Victor and Victorinus, who 
 were, with Claudian and his wife Bassa, subjected to many torments for a 
 period of three years, after which they were cast into prison, where they 
 ended their pilgrimage of life.
\switchcolumn*
\selectlanguage{latin}
In Cypro sancti Conónis Mártyris, qui sub Décio Imperatóre, clavis confíxus 
 pedes et ante currum jussus cúrrere, in génua procúbuit, atque in oratióne 
 réddidit spíritum.
\switchcolumn
\selectlanguage{english}
In Cyprus, in the time of Emperor Decius, St. Conon, martyr. He was 
 compelled to run before a chariot, with his feet pierced with nails, and 
 falling to his knees, he died in prayer.
\switchcolumn*
\selectlanguage{latin}
In Syria pássio sanctórum quadragínta duórum Mártyrum, qui, in Amório 
 comprehénsi et illuc perdúcti, ibi, egrégio perácto certámine, victóres 
 palmam martyrii percepérunt.
\switchcolumn
\selectlanguage{english}
In Syria, the passion of forty-two holy martyrs, who were arrested in 
 Amorium and taken to Syria, where they valiantly endured the test and 
 received the crown of martyrdom.
\switchcolumn*
\selectlanguage{latin}
Constantinópoli sancti Evágrii, qui, témpore Valéntis a Cathólicis eléctus 
 Epíscopus, et ab eódem Imperatóre in exsílium missus, Conféssor migrávit ad 
 Dóminum.v
\switchcolumn
\selectlanguage{english}
At Constantinople, St. Evagrius, who was elected Catholic bishop in the 
 reign of Valens, and being exiled by that emperor, later departed for 
 heaven.
\switchcolumn*
\selectlanguage{latin}
Bonóniæ sancti Basilíi Epíscopi, qui, a sancto 
 Silvéstro Papa ordinátus, verbo et exémplo créditam sibi Ecclésiam 
 sanctíssime gubernávit.
\switchcolumn
\selectlanguage{english}
At Bologna, St. Basil, bishop, who was ordained by Pope St. Sylvester, and 
 who governed the church entrusted to his care with great holiness, both by 
 word and example.
\switchcolumn*
\selectlanguage{latin}
Barcinóne, in Hispánia, beáti Ollegárii, primum Canónici, et póstea Epíscopi 
 Barcinonénsis, et Archiepíscopi Tarraconénsis.
\switchcolumn
\selectlanguage{english}
At Barcelona in Spain, blessed Ollegar, who was first a canon and afterwards 
 bishop of Barcelona and archbishop of Tarragona.
\switchcolumn*
\selectlanguage{latin}
Vitérbii beátæ Rosæ Vírginis, ex tértio Ordine 
 sancti Francísci.
\switchcolumn
\selectlanguage{english}
At Viterbo, blessed Rose, a virgin of the Third Order of St. Francis.
\switchcolumn*
\selectlanguage{latin}
Apud Gandávum, in Flándria, sanctæ Colétæ 
 Vírginis, quæ, primum tértii Ordinis Franciscális régulam proféssa, deínde, 
 divíno Spíritu affláta, quamplúra Moniálium secúndi ejúsdem Ordinis 
 monastéria primævæ restítuit disciplínæ; atque, divínis exornáta virtútibus 
 et innúmeris clara miráculis, a Pio Séptimo, Pontífice Máximo, in albo 
 Sanctórum adscrípta est.
\switchcolumn
\selectlanguage{english}
At Ghent in Flanders, St. Colette, virgin, who at first professed the rule 
 of the Third Order of St. Francis, and afterwards, by the inspiration of the 
 Holy Spirit, restored the pristine discipline to a great number of 
 monasteries of Nuns of the Second Order. Because she was graced 
 with heavenly virtues, and performed innumerable miracles, she was inscribed 
 on the roll of saints by Pope Pius VII.
\switchcolumn*
\selectlanguage{latin}
\end{paracol}


% ---- martyrology/mart03/mart0307.htm
\needspace{10\baselineskip}
\begin{paracol}{2}
\selectlanguage{latin}
\begin{center}{\color{gregoriocolor} Nonis Mártii. 
 Luna\dots\ }\end{center}
\switchcolumn
\selectlanguage{english}
\begin{center}{\color{gregoriocolor} The 7\textsuperscript{th} Day of March. 
 The\dots\ Day of the Moon.}\end{center}
\end{paracol}

\noindent\begin{tabularx}{\linewidth}{*{19}{>{\centering\arraybackslash}X}}
 \textcolor{gregoriocolor}{a} & \textcolor{gregoriocolor}{b} & \textcolor{gregoriocolor}{c} & \textcolor{gregoriocolor}{d} & \textcolor{gregoriocolor}{e} & \textcolor{gregoriocolor}{f} & \textcolor{gregoriocolor}{g} & \textcolor{gregoriocolor}{h} & \textcolor{gregoriocolor}{i} & \textcolor{gregoriocolor}{k} & \textcolor{gregoriocolor}{l} & \textcolor{gregoriocolor}{m} & \textcolor{gregoriocolor}{n} & \textcolor{gregoriocolor}{p} & \textcolor{gregoriocolor}{q} & \textcolor{gregoriocolor}{r} & \textcolor{gregoriocolor}{s} & \textcolor{gregoriocolor}{t} & \textcolor{gregoriocolor}{u} \\
 8 & 9 & 10 & 11 & 12 & 13 & 14 & 15 & 16 & 17 & 18 & 19 & 20 & 21 & 22 & 23 & 24 & 25 & 26 \\
\end{tabularx}
\vspace{0.5\baselineskip}
\noindent\begin{tabularx}{\linewidth}{*{12}{>{\centering\arraybackslash}X}}
 \textcolor{gregoriocolor}{A} & \textcolor{gregoriocolor}{B} & \textcolor{gregoriocolor}{C} & \textcolor{gregoriocolor}{D} & \textcolor{gregoriocolor}{E} & F & \textcolor{gregoriocolor}{F} & \textcolor{gregoriocolor}{G} & \textcolor{gregoriocolor}{H} & \textcolor{gregoriocolor}{M} & \textcolor{gregoriocolor}{N} & \textcolor{gregoriocolor}{P} \\
 27 & 28 & 29 & 30 & 1 & 2 & 2 & 3 & 4 & 5 & 6 & 7 \\
\end{tabularx}

\begin{paracol}{2}
\selectlanguage{latin}
\lettrine[lines=2]{I}{n} monastério Fossæ Novæ, prope Tarracínam, in 
 Campánia, sancti Thomæ Aquinátis, Confessóris et Ecclésiæ Doctóris, ex 
 Ordine Prædicatórum, nobilitáte géneris, vitæ sanctitáte et Theologíæ 
 sciéntia illustríssimi; quem Leo Papa Décimus tértius cæléstem Scholárum 
 ómnium catholicárum Patrónum declarávit.
\switchcolumn
\selectlanguage{english}
\lettrine[lines=2]{I}{n} the monastery of Fossanova, near Terracina in Campania, St. Thomas 
 Aquinas, confessor and doctor of the Church, a member of the Order of 
 Preachers, famous for his noble family, for the sanctity of his life, and 
 for his knowledge of theology. Pope Leo XIII declared him the heavenly 
 patron of all Catholic schools.
\switchcolumn*
\selectlanguage{latin}
Carthágine natális sanctárum Perpétuæ et 
 Felicitátis Mártyrum; e quibus Felícitas, cum esset prægnans (ut sanctus 
 Augustínus ait), juxta leges exspectáta ut páreret, dum parturiébat, dolébat, 
 objécta feris gaudébat. Passi quoque sunt cum eis Sátyrus, Saturnínus, 
 Revocátus et Secúndulus; quorum últimus quiévit in cárcere, réliqui omnes a 
 váriis béstiis sunt vexáti, ac demum gladiórum íctibus confécti, sub Sevéro 
 Príncipe, Sanctárum vero Perpétuæ et Felicitátis festum prídie hujus diéi 
 recólitur.
\switchcolumn
\selectlanguage{english}
At Carthage, the birthday of Saints Perpetua and Felicity, martyrs. 
 St. Augustine relates that Felicity being with child, her execution was 
 deferred, according to the law, until after her delivery. He states 
 that while she was in labour, she mourned, and when cast to the beasts, she 
 rejoiced. With them suffered Satyrus, Saturninus, Revocatus, and 
 Secundulus, the last of whom died in prison; the others were delivered to 
 the beasts, all during the reign of Severus. The feast of Saints 
 Perpetua and Felicity was celebrated yesterday.
\switchcolumn*
\selectlanguage{latin}
Cæsaréæ, in Palæstína, pássio sancti Eubúli, 
 qui fuit sócius sancti Hadriáni, atque, bíduo post illum, laniátus a 
 leónibus et gládio trucidátus, martyrii corónam, últimus ómnium in ea 
 civitáte, accépit.
\switchcolumn
\selectlanguage{english}
At Caesarea in Palestine, the passion of St. Eubulus, who was a companion of 
 St. Adrian. Two days after the latter's death, he was mangled by the 
 lions and put to death by the sword. He was the last of all those who 
 received the crown of martyrdom in that city.
\switchcolumn*
\selectlanguage{latin}
Nicomedíæ sancti Theóphili Epíscopi, qui, ob 
 cultum sanctárum Imáginum in exsílium pulsus, illic defúnctus est.
\switchcolumn
\selectlanguage{english}
At Nicomedia, St. Theophilus, bishop, who was driven into exile for the 
 veneration of sacred images, and died there.
\switchcolumn*
\selectlanguage{latin}
Pelúsii, in Ægypto, sancti Pauli Epíscopi, qui, 
 eándem ob causam, exsul occúbuit.
\switchcolumn
\selectlanguage{english}
At Pelusium in Egypt, St. Paul, bishop, who died in exile for the same 
 cause.
\switchcolumn*
\selectlanguage{latin}
Bríxiæ sancti Gaudiósi, Epíscopi et Confessóris.
\switchcolumn
\selectlanguage{english}
At Brescia, St. Gaudiosus, bishop and confessor.
\switchcolumn*
\selectlanguage{latin}
In Thebáide sancti Pauli, cognoménto Símplicis.
\switchcolumn
\selectlanguage{english}
In Thebais, St. Paul, surnamed the Simple.
\switchcolumn*
\selectlanguage{latin}
Floréntiæ, in Etrúria, sanctæ Terésiæ Margarítæ 
 Redi, Vírginis, Ordinis Carmelitárum Excalceatárum, vitæ puritáte ac 
 simplicitáte admirábilis, quam Pius Papa Undécimus sanctárum Vírginum albo 
 adscrípsit.
\switchcolumn
\selectlanguage{english}
At Florence in Etruria, St. Teresa Margaret Redi, virgin, a member of the 
 Order of Discalced Carmelites, of such admirable purity and simplicity that 
 Pope Pius XI solemnly enrolled her on the scroll of holy virgins.
\switchcolumn*
\selectlanguage{latin}
\end{paracol}


% ---- martyrology/mart03/mart0308.htm
\needspace{10\baselineskip}
\begin{paracol}{2}
\selectlanguage{latin}
\begin{center}{\color{gregoriocolor} Octávo Idus Mártii. 
 Luna\dots\ }\end{center}
\switchcolumn
\selectlanguage{english}
\begin{center}{\color{gregoriocolor} The 8\textsuperscript{th} Day of March. 
 The\dots\ Day of the Moon.}\end{center}
\end{paracol}

\noindent\begin{tabularx}{\linewidth}{*{19}{>{\centering\arraybackslash}X}}
 \textcolor{gregoriocolor}{a} & \textcolor{gregoriocolor}{b} & \textcolor{gregoriocolor}{c} & \textcolor{gregoriocolor}{d} & \textcolor{gregoriocolor}{e} & \textcolor{gregoriocolor}{f} & \textcolor{gregoriocolor}{g} & \textcolor{gregoriocolor}{h} & \textcolor{gregoriocolor}{i} & \textcolor{gregoriocolor}{k} & \textcolor{gregoriocolor}{l} & \textcolor{gregoriocolor}{m} & \textcolor{gregoriocolor}{n} & \textcolor{gregoriocolor}{p} & \textcolor{gregoriocolor}{q} & \textcolor{gregoriocolor}{r} & \textcolor{gregoriocolor}{s} & \textcolor{gregoriocolor}{t} & \textcolor{gregoriocolor}{u} \\
 9 & 10 & 11 & 12 & 13 & 14 & 15 & 16 & 17 & 18 & 19 & 20 & 21 & 22 & 23 & 24 & 25 & 26 & 27 \\
\end{tabularx}
\vspace{0.5\baselineskip}
\noindent\begin{tabularx}{\linewidth}{*{12}{>{\centering\arraybackslash}X}}
 \textcolor{gregoriocolor}{A} & \textcolor{gregoriocolor}{B} & \textcolor{gregoriocolor}{C} & \textcolor{gregoriocolor}{D} & \textcolor{gregoriocolor}{E} & F & \textcolor{gregoriocolor}{F} & \textcolor{gregoriocolor}{G} & \textcolor{gregoriocolor}{H} & \textcolor{gregoriocolor}{M} & \textcolor{gregoriocolor}{N} & \textcolor{gregoriocolor}{P} \\
 28 & 29 & 30 & 1 & 2 & 3 & 3 & 4 & 5 & 6 & 7 & 8 \\
\end{tabularx}

\begin{paracol}{2}
\selectlanguage{latin}
\lettrine[lines=2]{G}{ranátæ,} in Hispánia, sancti Joánnis de Deo 
 Confessóris, qui Ordinis Fratrum Hospitalitátis infirmórum fuit Institútor, 
 ac misericórdia in páuperes et sui despéctu éxstitit insígnis; quem Leo 
 Décimus tértius, Póntifex Máximus, cæléstem ómnium hospitálium et infirmórum 
 Patrónum renuntiávit.
\switchcolumn
\selectlanguage{english}
\lettrine[lines=2]{A}{t} Granada in Spain, St. John of God, founder of the Order of Brothers 
 Hospitallers, famed for his mercy to the poor, and his contempt of self. 
 Pope Leo XIII appointed him as heavenly patron of the sick and of all 
 hospitals.
\switchcolumn*
\selectlanguage{latin}
Nicomedíæ sancti Quinctílis, Epíscopi et 
 Mártyris.
\switchcolumn
\selectlanguage{english}
At Nicomedia, St. Quinctilis, bishop and martyr.
\switchcolumn*
\selectlanguage{latin}
In Africa sanctórum Mártyrum Cyrílli Epíscopi, Rogáti, Felícis, item Rogáti, 
 Beátæ, Heréniæ, Felicitátis, Urbáni, Silváni et 
 Mamílli.
\switchcolumn
\selectlanguage{english}
In Africa, the martyred Saints Cyril, bishop, Rogatus, Felix, another 
 Rogatus, Beata, Herenia, Felicitas, Urban, Silvanus, and Mamillus.
\switchcolumn*
\selectlanguage{latin}
Apud Antínoum, Ægypti urbem, natális sanctórum 
 Mártyrum Apollóni Diáconi, et Philémonis; qui, tenti et ad Júdicem addúcti, 
 et, cum sacrificáre idólis constánter renuíssent, ambo, perforátis calcáneis, 
 per civitátem horribíliter tracti, ac novíssime, gládio cæsi, martyrium 
 complevérunt.
\switchcolumn
\selectlanguage{english}
At Antinous, a city of Egypt, the birthday of the holy martyrs Apollonius, 
 deacon, and Philemon. They firmly refused to sacrifice to the idols, 
 and when arrested and brought to the judge they had their heels pierced, 
 were barbarously dragged through the city, at last completing their 
 martyrdom by being slain by the sword.
\switchcolumn*
\selectlanguage{latin}
Ibídem pássio sanctórum Ariáni Præsidis, 
 Theótici et aliórum trium; quos Judex submérsos in mare necávit, sed 
 delphinórum obséquio córpora eórum ad littus deláta sunt.
\switchcolumn
\selectlanguage{english}
In the same place, the passion of Saints Arian, governor, Theoticus, and 
 three others, whom the judge put to death by drowning in the sea. 
 Their bodies, however, were brought back by some dolphins.
\switchcolumn*
\selectlanguage{latin}
Carthágine sancti Póntii, qui fuit Diáconus beáti Cypriáni Epíscopi, et, 
 usque ad diem mortis illíus sústinens cum ipso exsílium, vitæ 
 et passiónis ejus egrégium volúmen relíquit, atque, in suis passiónibus 
 Dóminum semper gloríficans, corónam vitæ proméruit.
\switchcolumn
\selectlanguage{english}
At Carthage, St. Pontius, deacon of the blessed Cyprian, bishop, who 
 remained until death in exile with him, and composed an excellent history of 
 his life and martyrdom. By ever glorifying God in his own sufferings, 
 he merited the crown of life.
\switchcolumn*
\selectlanguage{latin}
Toléti, in Hispánia, deposítio beáti Juliáni, Epíscopi et Confessóris, 
 sanctitáte et doctrína celebérrimi.
\switchcolumn
\selectlanguage{english}
At Toledo in Spain, the death of blessed Julian, bishop and confessor, most 
 celebrated for his sanctity and learning.
\switchcolumn*
\selectlanguage{latin}
In Anglia sancti Felícis Epíscopi, qui orientáles Anglos ad fidem convértit.
\switchcolumn
\selectlanguage{english}
In England, St. Felix, bishop, who converted the East Angles to the faith.
\switchcolumn*
\selectlanguage{latin}
\end{paracol}


% ---- martyrology/mart03/mart0309.htm
\needspace{10\baselineskip}
\begin{paracol}{2}
\selectlanguage{latin}
\begin{center}{\color{gregoriocolor} Séptimo Idus Mártii. 
 Luna\dots\ }\end{center}
\switchcolumn
\selectlanguage{english}
\begin{center}{\color{gregoriocolor} The 9\textsuperscript{th} Day of March. 
 The\dots\ Day of the Moon.}\end{center}
\end{paracol}

\noindent\begin{tabularx}{\linewidth}{*{19}{>{\centering\arraybackslash}X}}
 \textcolor{gregoriocolor}{a} & \textcolor{gregoriocolor}{b} & \textcolor{gregoriocolor}{c} & \textcolor{gregoriocolor}{d} & \textcolor{gregoriocolor}{e} & \textcolor{gregoriocolor}{f} & \textcolor{gregoriocolor}{g} & \textcolor{gregoriocolor}{h} & \textcolor{gregoriocolor}{i} & \textcolor{gregoriocolor}{k} & \textcolor{gregoriocolor}{l} & \textcolor{gregoriocolor}{m} & \textcolor{gregoriocolor}{n} & \textcolor{gregoriocolor}{p} & \textcolor{gregoriocolor}{q} & \textcolor{gregoriocolor}{r} & \textcolor{gregoriocolor}{s} & \textcolor{gregoriocolor}{t} & \textcolor{gregoriocolor}{u} \\
 10 & 11 & 12 & 13 & 14 & 15 & 16 & 17 & 18 & 19 & 20 & 21 & 22 & 23 & 24 & 25 & 26 & 27 & 28 \\
\end{tabularx}
\vspace{0.5\baselineskip}
\noindent\begin{tabularx}{\linewidth}{*{12}{>{\centering\arraybackslash}X}}
 \textcolor{gregoriocolor}{A} & \textcolor{gregoriocolor}{B} & \textcolor{gregoriocolor}{C} & \textcolor{gregoriocolor}{D} & \textcolor{gregoriocolor}{E} & F & \textcolor{gregoriocolor}{F} & \textcolor{gregoriocolor}{G} & \textcolor{gregoriocolor}{H} & \textcolor{gregoriocolor}{M} & \textcolor{gregoriocolor}{N} & \textcolor{gregoriocolor}{P} \\
 29 & 30 & 1 & 2 & 3 & 4 & 4 & 5 & 6 & 7 & 8 & 9 \\
\end{tabularx}

\begin{paracol}{2}
\selectlanguage{latin}
\lettrine[lines=2]{R}{omæ} sanctæ Francíscæ Víduæ, nobilitáte géneris, 
 vitæ sanctitáte et miraculórum dono célebris.
\switchcolumn
\selectlanguage{english}
\lettrine[lines=2]{A}{t} Rome, St. Frances, widow, renowned for her noble family, holy life, and 
 the gift of miracles.
\switchcolumn*
\selectlanguage{latin}
Apud Sebásten, in Arménia, natális sanctórum Quadragínta mílitum Cappadócum, 
 qui, témpore Licínii Imperatóris, sub Præside 
 Agricoláo, post víncula et cárceres tetérrimos, post cæsas lapídibus fácies, 
 nudi sub dio, frigidíssimo híemis témpore, supra stagnum rigens pernoctáre 
 jussi sunt, ubi gelu constrícta eórum córpora disrumpebántur, ac demum 
 crurifrágio martyrium consummárunt. Erant autem inter eos nobilióres 
 Cyrion et Cándidus; eorúmque ómnium præcláras glórias sanctus Basilíus 
 aliíque Patres scriptis suis celebrárunt. Ipsórum porro Mártyrum 
 festívitas sequénti die recólitur.
\switchcolumn
\selectlanguage{english}
At Sebaste in Armenia, under the governor Agricolaus, in the time of Emperor 
 Licinius, the birthday of forty holy soldiers of Cappadocia. After 
 being chained down in foul dungeons, after having their faces bruised with 
 stones, and being condemned to spend the night naked, in the open during the 
 coldest part of winter, on a frozen lake where their bodies were benumbed 
 and covered with ice, they completed their martyrdom by having their limbs 
 crushed. The most noteworthy among them were Cyrion and Candidus. 
 Their glorious triumph has been celebrated by St. Basil and other Fathers in 
 their writings. Their feast is kept tomorrow.
\switchcolumn*
\selectlanguage{latin}
Nyssæ deposítio sancti Gregórii Epíscopi, qui 
 sanctórum Basilíi et Emméliæ fílius, et sanctórum item Basilíi Magni ac 
 Petri Sebasténsis Episcopórum et Macrínæ Vírginis frater éxstitit; atque, 
 vita et eruditióne claríssimus, ob fídei cathólicæ defensiónem, sub Ariáno 
 Imperatóre Valénte, civitáte sua pulsus est.
\switchcolumn
\selectlanguage{english}
At Nyssa, the death of St. Gregory, the son of Saints Basil and Emmelia, and 
 the brother of Saints Basil the Great, bishop, and Peter, bishop of Sebaste, 
 and Macrina, virgin. His life and his great learning brought him fame. 
 He was driven from his own city for having defended the Catholic faith 
 during the reign of the Arian emperor Valens.
\switchcolumn*
\selectlanguage{latin}
Barcinóne, in Hispánia, sancti Paciáni Epíscopi, tam vita quam sermóne 
 conspícui; qui, témpore Theodósii Príncipis, in 
 última senectúte finem vitæ 
 sortítus est.
\switchcolumn
\selectlanguage{english}
At Barcelona in Spain, Bishop St. Pacian, distinguished by his life and 
 preaching. He ended his career in extreme old age, in the time of 
 Emperor Theodosius.
\switchcolumn*
\selectlanguage{latin}
Bonóniæ sanctæ Catharínæ Vírginis, e secúndo 
 Ordine sancti Francísci, quæ vitæ sanctitáte fuit illústris. Ipsíus 
 autem corpus magno cum honóre ibídem cólitur.
\switchcolumn
\selectlanguage{english}
At Bologna, St. Catherine, virgin, of the Second Order of St. Francis, 
 illustrious for the holiness of her life. Her body is greatly honoured 
 in that city.
\switchcolumn*
\selectlanguage{latin}
\end{paracol}


% ---- martyrology/mart03/mart0310.htm
\needspace{10\baselineskip}
\begin{paracol}{2}
\selectlanguage{latin}
\begin{center}{\color{gregoriocolor} Sexto Idus Mártii. 
 Luna\dots\ }\end{center}
\switchcolumn
\selectlanguage{english}
\begin{center}{\color{gregoriocolor} The 10\textsuperscript{th} Day of March. 
 The\dots\ Day of the Moon.}\end{center}
\end{paracol}

\noindent\begin{tabularx}{\linewidth}{*{19}{>{\centering\arraybackslash}X}}
 \textcolor{gregoriocolor}{a} & \textcolor{gregoriocolor}{b} & \textcolor{gregoriocolor}{c} & \textcolor{gregoriocolor}{d} & \textcolor{gregoriocolor}{e} & \textcolor{gregoriocolor}{f} & \textcolor{gregoriocolor}{g} & \textcolor{gregoriocolor}{h} & \textcolor{gregoriocolor}{i} & \textcolor{gregoriocolor}{k} & \textcolor{gregoriocolor}{l} & \textcolor{gregoriocolor}{m} & \textcolor{gregoriocolor}{n} & \textcolor{gregoriocolor}{p} & \textcolor{gregoriocolor}{q} & \textcolor{gregoriocolor}{r} & \textcolor{gregoriocolor}{s} & \textcolor{gregoriocolor}{t} & \textcolor{gregoriocolor}{u} \\
 11 & 12 & 13 & 14 & 15 & 16 & 17 & 18 & 19 & 20 & 21 & 22 & 23 & 24 & 25 & 26 & 27 & 28 & 29 \\
\end{tabularx}
\vspace{0.5\baselineskip}
\noindent\begin{tabularx}{\linewidth}{*{12}{>{\centering\arraybackslash}X}}
 \textcolor{gregoriocolor}{A} & \textcolor{gregoriocolor}{B} & \textcolor{gregoriocolor}{C} & \textcolor{gregoriocolor}{D} & \textcolor{gregoriocolor}{E} & F & \textcolor{gregoriocolor}{F} & \textcolor{gregoriocolor}{G} & \textcolor{gregoriocolor}{H} & \textcolor{gregoriocolor}{M} & \textcolor{gregoriocolor}{N} & \textcolor{gregoriocolor}{P} \\
 30 & 1 & 2 & 3 & 4 & 5 & 5 & 6 & 7 & 8 & 9 & 10 \\
\end{tabularx}

\begin{paracol}{2}
\selectlanguage{latin}
\lettrine[lines=2]{S}{anctórum} Quadragínta Mártyrum, quorum natális prídie hujus diéi recensétur.
\switchcolumn
\selectlanguage{english}
\lettrine[lines=2]{T}{he} forty holy martyrs whose birthday was commemorated yesterday.
\switchcolumn*
\selectlanguage{latin}
Apaméæ, in Phrygia, natális sanctórum Mártyrum 
 Caji et Alexándri, qui (ut scribit Apollináris, Hierapolitánus Epíscopus, in 
 libro advérsus Cataphrygas hæréticos), in persecutióne Marci Antoníni et 
 Lúcii Veri, glorióso martyrio coronáti sunt.
\switchcolumn
\selectlanguage{english}
At Apamea in Phrygia, during the persecution of Marcus Antoninus and Lucius 
 Verus, the birthday of the holy martyrs Caius and Alexander. They were 
 crowned with a glorious martyrdom, as is related by Apollinaris, bishop of 
 Hierapolis, in his book against the Cataphrygian heretics.
\switchcolumn*
\selectlanguage{latin}
In Pérside pássio sanctórum quadragínta duórum Mártyrum.
\switchcolumn
\selectlanguage{english}
In Persia, the passion of forty-two holy martyrs.
\switchcolumn*
\selectlanguage{latin}
Corínthi sanctórum Mártyrum Codráti, Dionysii, Cypriáni, Anécti, Pauli et 
 Crescéntis, qui, in persecutióne Décii et Valeriáni, sub Jásone 
 Prǽside, gládio cæsi sunt.
\switchcolumn
\selectlanguage{english}
At Corinth, the holy martyrs Codratus, Denis, Cyprian, Anectus, Paul, and 
 Crescens, who were slain with the sword in the persecution of Decius and 
 Valerian, under Jason, the governor.
\switchcolumn*
\selectlanguage{latin}
In Africa sancti Victóris Mártyris, in cujus solemnitáte sanctus Augustínus 
 ad pópulum de ipso tractátum hábuit.
\switchcolumn
\selectlanguage{english}
In Africa, St. Victor, martyr, on whose feast day St. Augustine delivered a 
 sermon to his people.
\switchcolumn*
\selectlanguage{latin}
Romæ sancti Simplícii, Papæ et Confessóris.
\switchcolumn
\selectlanguage{english}
At Rome, St. Simplicius, pope and confessor.
\switchcolumn*
\selectlanguage{latin}
Hierosólymis sancti Macárii, Epíscopi et Confessóris; cujus hortátu loca 
 sancta a Constantíno Magno et beáta Hélena, 
 ejus matre, expurgáta sunt et sacris Basílicis illustráta.
\switchcolumn
\selectlanguage{english}
At Jerusalem, St. Macarius, bishop and confessor, at whose exhortation the 
 holy places were purged by Constantine the Great and St. Helen, his mother, 
 and beautified by sacred basilicas.
\switchcolumn*
\selectlanguage{latin}
Lutétiæ Parisiórum deposítio sancti Droctovéi 
 Abbátis, qui fuit discípulus beáti Germáni Epíscopi.
\switchcolumn
\selectlanguage{english}
At Paris, the death of Abbot St. Droctoveus, who was a disciple of the 
 saintly Bishop Germanus.
\switchcolumn*
\selectlanguage{latin}
In monastério Bobiénsi sancti Attalæ Abbátis, 
 miráculis clari.
\switchcolumn
\selectlanguage{english}
In the monastery of Bobbio, St. Attala, abbot, renowned for his miracles.
\switchcolumn*
\selectlanguage{latin}
\end{paracol}


% ---- martyrology/mart03/mart0311.htm
\needspace{10\baselineskip}
\begin{paracol}{2}
\selectlanguage{latin}
\begin{center}{\color{gregoriocolor} Quinto Idus Mártii. 
 Luna\dots\ }\end{center}
\switchcolumn
\selectlanguage{english}
\begin{center}{\color{gregoriocolor} The 11\textsuperscript{th} Day of March. 
 The\dots\ Day of the Moon.}\end{center}
\end{paracol}

\noindent\begin{tabularx}{\linewidth}{*{19}{>{\centering\arraybackslash}X}}
 \textcolor{gregoriocolor}{a} & \textcolor{gregoriocolor}{b} & \textcolor{gregoriocolor}{c} & \textcolor{gregoriocolor}{d} & \textcolor{gregoriocolor}{e} & \textcolor{gregoriocolor}{f} & \textcolor{gregoriocolor}{g} & \textcolor{gregoriocolor}{h} & \textcolor{gregoriocolor}{i} & \textcolor{gregoriocolor}{k} & \textcolor{gregoriocolor}{l} & \textcolor{gregoriocolor}{m} & \textcolor{gregoriocolor}{n} & \textcolor{gregoriocolor}{p} & \textcolor{gregoriocolor}{q} & \textcolor{gregoriocolor}{r} & \textcolor{gregoriocolor}{s} & \textcolor{gregoriocolor}{t} & \textcolor{gregoriocolor}{u} \\
 12 & 13 & 14 & 15 & 16 & 17 & 18 & 19 & 20 & 21 & 22 & 23 & 24 & 25 & 26 & 27 & 28 & 29 & 30 \\
\end{tabularx}
\vspace{0.5\baselineskip}
\noindent\begin{tabularx}{\linewidth}{*{12}{>{\centering\arraybackslash}X}}
 \textcolor{gregoriocolor}{A} & \textcolor{gregoriocolor}{B} & \textcolor{gregoriocolor}{C} & \textcolor{gregoriocolor}{D} & \textcolor{gregoriocolor}{E} & F & \textcolor{gregoriocolor}{F} & \textcolor{gregoriocolor}{G} & \textcolor{gregoriocolor}{H} & \textcolor{gregoriocolor}{M} & \textcolor{gregoriocolor}{N} & \textcolor{gregoriocolor}{P} \\
 1 & 2 & 3 & 4 & 5 & 6 & 6 & 7 & 8 & 9 & 10 & 11 \\
\end{tabularx}

\begin{paracol}{2}
\selectlanguage{latin}
\lettrine[lines=2]{S}{ardis} sancti Euthymii Epíscopi, qui, ob cultum sanctárum Imáginum, a 
 Michaéle, Imperatóre Iconoclásta, in exsílium missus est; ac demum, 
 Theóphilo imperánte, búbulis nervis, 
 inhumániter cæsus, martyrium consummávit.
\switchcolumn
\selectlanguage{english}
\lettrine[lines=2]{A}{t} Sardis, St. Euthymius, bishop, who was sent into exile by the Iconoclast 
 emperor Michael for the veneration of sacred images. Afterwards, in 
 the reign of Theophilus, he was barbarously beaten with knotted clubs, which 
 completed his martyrdom.
\switchcolumn*
\selectlanguage{latin}
Córdubæ, in Hispánia, sancti Eulógii, 
 Presbyteri et Mártyris; qui, in persecutióne Saracenórum, ob præcláram et 
 intrépidam Christi confessiónem, verbéribus et álapis cæsus ac decollátus 
 gládio, adjúngi ejúsdem urbis Martyribus méruit, quorum pro fide certámina 
 scribéndo fúerat æmulátus.
\switchcolumn
\selectlanguage{english}
At Cordova in Spain, St. Eulogius, priest, who deserved to be associated 
 with the martyrs of that city because, in writing of their trials for the 
 faith, he had envied their happiness. On account of his own fearless 
 and intrepid confession of Christ, he was scourged and beaten with rods, and 
 finally beheaded during the Saracen persecution.
\switchcolumn*
\selectlanguage{latin}
Carthágine sanctórum Mártyrum Heráclíi et 
 Zósimi.
\switchcolumn
\selectlanguage{english}
At Carthage, the holy martyrs Heraclius and Zosimus.
\switchcolumn*
\selectlanguage{latin}
Alexandríæ pássio sanctórum Cándidi, Piperiónis 
 et aliórum vigínti.
\switchcolumn
\selectlanguage{english}
At Alexandria, the passion of Saints Candidus, Piperion, and twenty others.
\switchcolumn*
\selectlanguage{latin}
Laodicéæ, in Syria, sanctórum Mártyrum Tróphimi 
 et Thali, qui in persecutióne Diocletiáni, post multa sǽvaque torménta, 
 corónas glóriæ sunt assecúti.
\switchcolumn
\selectlanguage{english}
At Laodicea in Syria, during the persecution of Diocletian, the holy martyrs 
 Trophimus and Thalus, who obtained their crowns of glory after many severe 
 torments.
\switchcolumn*
\selectlanguage{latin}
Antiochíæ commemorátio plurimórum sanctórum 
 Mártyrum, quorum álii, Maximiáni Imperatóris mandáto, candéntibus cratículis 
 superpósiti, non ad mortem sed a diutúrnum cruciátum assáti, álii áliis 
 sævíssimis affécti supplíciis, ad palmam martyrii pervenérunt.
\switchcolumn
\selectlanguage{english}
At Antioch, the Commemoration of many holy martyrs, some of whom by order of 
 Emperor Maximian were laid on red hot gridirons, not to be burned to death, 
 but to continue their suffering a longer time; others were subjected to 
 different horrible torments, and won the palm of martyrdom.
\switchcolumn*
\selectlanguage{latin}
Item sanctórum Mártyrum Gorgónii et Firmi.
\switchcolumn
\selectlanguage{english}
Also, Saints Gorgonius and Firmus.
\switchcolumn*
\selectlanguage{latin}
Hierosólymis sancti Sophrónii Epíscopi.
\switchcolumn
\selectlanguage{english}
At Jerusalem, Bishop St. Sophronius.
\switchcolumn*
\selectlanguage{latin}
Medioláni sancti Benedícti Epíscopi.
\switchcolumn
\selectlanguage{english}
At Milan, St. Benedict, bishop.
\switchcolumn*
\selectlanguage{latin}
In fínibus Ambianénsium sancti Firmíni Abbátis.
\switchcolumn
\selectlanguage{english}
In the diocese of Amiens, St. Firmin, abbot.
\switchcolumn*
\selectlanguage{latin}
Carthágine sancti Constantíni Confessóris.
\switchcolumn
\selectlanguage{english}
At Carthage, St. Constantine, confessor.
\switchcolumn*
\selectlanguage{latin}
Babúci, in Hérnicis, sancti Petri Confessóris, 
 miraculórum glória conspícui.
\switchcolumn
\selectlanguage{english}
At Babucum in Campania, St. Peter, confessor, who was renowned for his 
 miracles.
\switchcolumn*
\selectlanguage{latin}
\end{paracol}


% ---- martyrology/mart03/mart0312.htm
\needspace{10\baselineskip}
\begin{paracol}{2}
\selectlanguage{latin}
\begin{center}{\color{gregoriocolor} Quarto Idus Mártii. 
 Luna\dots\ }\end{center}
\switchcolumn
\selectlanguage{english}
\begin{center}{\color{gregoriocolor} The 12\textsuperscript{th} Day of March. 
 The\dots\ Day of the Moon.}\end{center}
\end{paracol}

\noindent\begin{tabularx}{\linewidth}{*{19}{>{\centering\arraybackslash}X}}
 \textcolor{gregoriocolor}{a} & \textcolor{gregoriocolor}{b} & \textcolor{gregoriocolor}{c} & \textcolor{gregoriocolor}{d} & \textcolor{gregoriocolor}{e} & \textcolor{gregoriocolor}{f} & \textcolor{gregoriocolor}{g} & \textcolor{gregoriocolor}{h} & \textcolor{gregoriocolor}{i} & \textcolor{gregoriocolor}{k} & \textcolor{gregoriocolor}{l} & \textcolor{gregoriocolor}{m} & \textcolor{gregoriocolor}{n} & \textcolor{gregoriocolor}{p} & \textcolor{gregoriocolor}{q} & \textcolor{gregoriocolor}{r} & \textcolor{gregoriocolor}{s} & \textcolor{gregoriocolor}{t} & \textcolor{gregoriocolor}{u} \\
 13 & 14 & 15 & 16 & 17 & 18 & 19 & 20 & 21 & 22 & 23 & 24 & 25 & 26 & 27 & 28 & 29 & 30 & 1 \\
\end{tabularx}
\vspace{0.5\baselineskip}
\noindent\begin{tabularx}{\linewidth}{*{12}{>{\centering\arraybackslash}X}}
 \textcolor{gregoriocolor}{A} & \textcolor{gregoriocolor}{B} & \textcolor{gregoriocolor}{C} & \textcolor{gregoriocolor}{D} & \textcolor{gregoriocolor}{E} & F & \textcolor{gregoriocolor}{F} & \textcolor{gregoriocolor}{G} & \textcolor{gregoriocolor}{H} & \textcolor{gregoriocolor}{M} & \textcolor{gregoriocolor}{N} & \textcolor{gregoriocolor}{P} \\
 2 & 3 & 4 & 5 & 6 & 7 & 7 & 8 & 9 & 10 & 11 & 12 \\
\end{tabularx}

\begin{paracol}{2}
\selectlanguage{latin}
\lettrine[lines=2]{R}{omæ} sancti Gregórii Primi, Papæ, Confessóris 
 et Ecclésiæ Doctóris exímii; qui, ob res præcláræ gestas atque Anglos ad 
 Christi fidem convérsos, Magnus est dictus et Anglórum Apóstolus appellátus.
\switchcolumn
\selectlanguage{english}
\lettrine[lines=2]{A}{t} Rome, St. Gregory, pope and eminent doctor of the Church, who on account 
 of his illustrious deeds and the conversion of the English to the faith of 
 Christ, was surnamed the Great, and called the Apostle of England.
\switchcolumn*
\selectlanguage{latin}
Ibídem deposítio sancti Innocéntii Primi, Papæ 
 et Confessóris. Ipsíus autem festum quinto Kaléndas Augústi celebrátur.
\switchcolumn
\selectlanguage{english}
In the same place, the death of St. Innocent I, pope and confessor. 
 His feast is celebrated on the 28th of July.
\switchcolumn*
\selectlanguage{latin}
Item Romæ sancti Mamiliáni Mártyris.
\switchcolumn
\selectlanguage{english}
Also at Rome, St. Mamilian, martyr.
\switchcolumn*
\selectlanguage{latin}
Nicomedíæ sanctórum Egdúni Presbyteri, et 
 aliórum septem; qui singuli diébus suffocáti sunt, ut céteris metus 
 incuterétur.
\switchcolumn
\selectlanguage{english}
At Nicomedia, St. Egdunus, priest, and seven others, who, one by one, on 
 successive days, were strangled in order to terrify those who remained.
\switchcolumn*
\selectlanguage{latin}
Ibídem pássio sancti Petri Mártyris, qui, cum esset cubiculárius Diocletiáni 
 Imperatóris, et libérius de imménsis Mártyrum supplíciis quererétur, 
 proptérea, jubénte eódem, in médium addúcitur, ac primo suspénsus, 
 diutíssime flagris torquétur, deínde acéto ac sale perfúsus, ad últimum in 
 cratícula lento igne assátur, sicque vere Petri éxstitit et fídei heres et 
 nóminis.
\switchcolumn
\selectlanguage{english}
In the same city, the passion of the martyr St. Peter, chamberlain to 
 Emperor Diocletian. For openly complaining of the atrocious torments 
 inflicted upon the martyrs, he was, by order of the emperor, first suspended 
 and for a long time scourged, then had salt and vinegar poured on his 
 wounds, and finally was burned on a grate over a slow fire. Thus did 
 he become a true heir of St. Peter's name and faith.
\switchcolumn*
\selectlanguage{latin}
Constantinópoli sancti Theóphanis, qui, ex ditíssimo pauper Mónachus 
 efféctus, ab ímpio Leóne Arméno, pro cultu sacrárum Imáginum, biénnio 
 deténtus est in cárcere, et inde in Samothráciam deportátus, ibídem,
 ærúmnis conféctus, réddidit spíritum, multísque 
 miráculis cláruit.
\switchcolumn
\selectlanguage{english}
At Constantinople, St. Theophanes, who gave up great riches to embrace the 
 poverty of the monastic state. The impious Leo the Armenian kept him 
 in prison for two years because of his veneration of sacred images, and 
 later sent him into Thrace in exile. There, overwhelmed with 
 afflictions, but famous for miracles, death came upon him.
\switchcolumn*
\selectlanguage{latin}
Cápuæ sancti Bernárdi, Epíscopi et Confessóris.
\switchcolumn
\selectlanguage{english}
At Capua, St. Bernard, bishop and confessor.
\switchcolumn*
\selectlanguage{latin}
\end{paracol}


% ---- martyrology/mart03/mart0313.htm
\needspace{10\baselineskip}
\begin{paracol}{2}
\selectlanguage{latin}
\begin{center}{\color{gregoriocolor} Tértio Idus Mártii. 
 Luna\dots\ }\end{center}
\switchcolumn
\selectlanguage{english}
\begin{center}{\color{gregoriocolor} The 13\textsuperscript{th} Day of March. 
 The\dots\ Day of the Moon.}\end{center}
\end{paracol}

\noindent\begin{tabularx}{\linewidth}{*{19}{>{\centering\arraybackslash}X}}
 \textcolor{gregoriocolor}{a} & \textcolor{gregoriocolor}{b} & \textcolor{gregoriocolor}{c} & \textcolor{gregoriocolor}{d} & \textcolor{gregoriocolor}{e} & \textcolor{gregoriocolor}{f} & \textcolor{gregoriocolor}{g} & \textcolor{gregoriocolor}{h} & \textcolor{gregoriocolor}{i} & \textcolor{gregoriocolor}{k} & \textcolor{gregoriocolor}{l} & \textcolor{gregoriocolor}{m} & \textcolor{gregoriocolor}{n} & \textcolor{gregoriocolor}{p} & \textcolor{gregoriocolor}{q} & \textcolor{gregoriocolor}{r} & \textcolor{gregoriocolor}{s} & \textcolor{gregoriocolor}{t} & \textcolor{gregoriocolor}{u} \\
 14 & 15 & 16 & 17 & 18 & 19 & 20 & 21 & 22 & 23 & 24 & 25 & 26 & 27 & 28 & 29 & 30 & 1 & 2 \\
\end{tabularx}
\vspace{0.5\baselineskip}
\noindent\begin{tabularx}{\linewidth}{*{12}{>{\centering\arraybackslash}X}}
 \textcolor{gregoriocolor}{A} & \textcolor{gregoriocolor}{B} & \textcolor{gregoriocolor}{C} & \textcolor{gregoriocolor}{D} & \textcolor{gregoriocolor}{E} & F & \textcolor{gregoriocolor}{F} & \textcolor{gregoriocolor}{G} & \textcolor{gregoriocolor}{H} & \textcolor{gregoriocolor}{M} & \textcolor{gregoriocolor}{N} & \textcolor{gregoriocolor}{P} \\
 3 & 4 & 5 & 6 & 7 & 8 & 8 & 9 & 10 & 11 & 12 & 13 \\
\end{tabularx}

\begin{paracol}{2}
\selectlanguage{latin}
\lettrine[lines=2]{C}{órdubæ} in Hispánia, sanctórum Mártyrum 
 Ruderíci Presbyteri, et Salomónis.
\switchcolumn
\selectlanguage{english}
\lettrine[lines=2]{A}{t} Cordova in Spain, the holy martyrs Roderick, priest, and Solomon.
\switchcolumn*
\selectlanguage{latin}
Nicomediæ natális sanctórum Mártyrum Macedónii, 
 Patríciæ uxóris, et Modéstæ fíliæ.
\switchcolumn
\selectlanguage{english}
At Nicomedia, the birthday of the holy martyrs Macedonius, Patricia, his 
 wife, and his daughter Modesta.
\switchcolumn*
\selectlanguage{latin}
Nicǽæ, in Bithynia, sanctórum Mártyrum Theusétæ, 
 ejúsque fílii Horris, Theodóræ, Nymphodóræ, Marci et Arábiæ; qui omnes pro 
 Christo igni tráditi sunt.
\switchcolumn
\selectlanguage{english}
At Nicaea in Bithynia, the holy martyrs Theusetas and Horres, his son; 
 Theodore, Nymphodora, Mark, and Arabia, who were all burned to death for 
 Christ.
\switchcolumn*
\selectlanguage{latin}
Hermópoli, in Ægypto, sancti Sabíni Mártyris, 
 qui multa passus, tandem, projéctus in flumen, martyrium consummávit.
\switchcolumn
\selectlanguage{english}
At Hermopolis in Egypt, the martyr St. Sabinus, who suffered many torments, 
 and at last completed his martyrdom by being cast into a river.
\switchcolumn*
\selectlanguage{latin}
In Pérside sanctæ Christínæ, Vírginis et 
 Mártyris.
\switchcolumn
\selectlanguage{english}
In Persia, St. Christina, virgin and martyr.
\switchcolumn*
\selectlanguage{latin}
Apud Camerínum sancti Ansovíni, Epíscopi et Confessóris.
\switchcolumn
\selectlanguage{english}
At Camerino, St. Ansovinus, bishop and confessor.
\switchcolumn*
\selectlanguage{latin}
In Thebáide deposítio sanctæ Euphrásiæ Vírginis.
\switchcolumn
\selectlanguage{english}
In Thebais, the death of St. Euphrasia, virgin.
\switchcolumn*
\selectlanguage{latin}
Constantinópoli Translátio sancti Nicéphori, 
 Epíscopi ejúsdem urbis et Confessóris; cujus corpus e Proconnéso, 
 Propóntidis ínsula, ubi ipse quarto Nonas Júnii ob sanctárum Imáginum cultum 
 exsul obíerat, Constantinópolim relátum est, atque a sancto illíus civitátis 
 Epíscopo Methódio honorífice in templo sanctórum Apostolórum sepúltum, hac 
 ipsa recurrénte die, in qua olim idem Nicéphorus in exsílium fúerat 
 deportátus.
\switchcolumn
\selectlanguage{english}
At Constantinople, the transferral of the body of St. Nicephorus, bishop of 
 that city, and confessor. The body was returned from the island of 
 Propontis in the Proconnesus, where his death occurred on the 5th of June 
 while in exile for his reverence of sacred images. He was buried with 
 honour by Bishop Methodius in the Church of the Holy Apostles on this the 
 anniversary day of his exile.
\switchcolumn*
\selectlanguage{latin}
\end{paracol}


% ---- martyrology/mart03/mart0314.htm
\needspace{10\baselineskip}
\begin{paracol}{2}
\selectlanguage{latin}
\begin{center}{\color{gregoriocolor} Prídie Idus Mártii. 
 Luna\dots\ }\end{center}
\switchcolumn
\selectlanguage{english}
\begin{center}{\color{gregoriocolor} The 14\textsuperscript{th} Day of March. 
 The\dots\ Day of the Moon.}\end{center}
\end{paracol}

\noindent\begin{tabularx}{\linewidth}{*{19}{>{\centering\arraybackslash}X}}
 \textcolor{gregoriocolor}{a} & \textcolor{gregoriocolor}{b} & \textcolor{gregoriocolor}{c} & \textcolor{gregoriocolor}{d} & \textcolor{gregoriocolor}{e} & \textcolor{gregoriocolor}{f} & \textcolor{gregoriocolor}{g} & \textcolor{gregoriocolor}{h} & \textcolor{gregoriocolor}{i} & \textcolor{gregoriocolor}{k} & \textcolor{gregoriocolor}{l} & \textcolor{gregoriocolor}{m} & \textcolor{gregoriocolor}{n} & \textcolor{gregoriocolor}{p} & \textcolor{gregoriocolor}{q} & \textcolor{gregoriocolor}{r} & \textcolor{gregoriocolor}{s} & \textcolor{gregoriocolor}{t} & \textcolor{gregoriocolor}{u} \\
 15 & 16 & 17 & 18 & 19 & 20 & 21 & 22 & 23 & 24 & 25 & 26 & 27 & 28 & 29 & 30 & 1 & 2 & 3 \\
\end{tabularx}
\vspace{0.5\baselineskip}
\noindent\begin{tabularx}{\linewidth}{*{12}{>{\centering\arraybackslash}X}}
 \textcolor{gregoriocolor}{A} & \textcolor{gregoriocolor}{B} & \textcolor{gregoriocolor}{C} & \textcolor{gregoriocolor}{D} & \textcolor{gregoriocolor}{E} & F & \textcolor{gregoriocolor}{F} & \textcolor{gregoriocolor}{G} & \textcolor{gregoriocolor}{H} & \textcolor{gregoriocolor}{M} & \textcolor{gregoriocolor}{N} & \textcolor{gregoriocolor}{P} \\
 4 & 5 & 6 & 7 & 8 & 9 & 9 & 10 & 11 & 12 & 13 & 14 \\
\end{tabularx}

\begin{paracol}{2}
\selectlanguage{latin}
\lettrine[lines=2]{R}{omæ,} in agro Veráno, sancti Leónis, Epíscopi 
 et Mártyris.
\switchcolumn
\selectlanguage{english}
\lettrine[lines=2]{A}{t} Rome, in the Veranian Field, St. Leo, bishop and martyr.
\switchcolumn*
\selectlanguage{latin}
Item Romæ natális sanctórum quadragínta septem 
 Mártyrum, qui baptizáti sunt a beáto Apóstolo Petro, cum tenerétur in 
 custódia Mamertíni cum Coapóstolo suo Paulo, ubi novem menses deténti sunt; 
 qui omnes, sub devotíssima fídei confessióne, Neroniáno gládio consúmpti 
 sunt.
\switchcolumn
\selectlanguage{english}
Also at Rome, the birthday of forty-seven holy martyrs who were baptized by 
 the apostle St. Peter while in the Mamertine Prison with St. Paul his fellow 
 apostle. After an imprisonment of nine months, they all fell by the 
 sword of Nero for their generous confession of faith.
\switchcolumn*
\selectlanguage{latin}
In Província Valériæ sanctórum duórum 
 Monachórum, quos Longobárdi suspéndio necavérunt in árbore; in qua Mártyres, 
 licet defúncti, ab hóstibus ipsis audíti sunt psállere.
\switchcolumn
\selectlanguage{english}
In the province of Valeria, two saintly monks, who were hanged on a tree by 
 the Lombards, and although dead, were heard singing psalms even by their 
 enemies.
\switchcolumn*
\selectlanguage{latin}
In ea étiam persecutióne Diáconus Ecclésiæ 
 Marsicánæ, in confessióne fídei, cápite truncátus est.
\switchcolumn
\selectlanguage{english}
In the same persecution, a deacon of the church of Marsico who was beheaded 
 for the confession of faith.
\switchcolumn*
\selectlanguage{latin}
In Africa sanctórum Mártyrum Petri et Aphrodísii, qui, in persecutióne 
 Wandálica, martyrii corónam percepérunt.
\switchcolumn
\selectlanguage{english}
In Africa, the holy martyrs Peter and Aphrodisius, who received the crown of 
 martyrdom in the Vandal persecution.
\switchcolumn*
\selectlanguage{latin}
Carrhis, in Mesopotámia, sancti Eutychii patrícii, et Sociórum, qui ab 
 Evelid, Arabum Rege, ob fídei confessiónem, interémpti sunt.
\switchcolumn
\selectlanguage{english}
At Carrhae in Mesopotamia, the patrician St. Eutychius and his companions, 
 who were killed by Evelid, king of Arabia, for the confession of the faith.
\switchcolumn*
\selectlanguage{latin}
Halberstátti, in Germánia, dormítio beátæ 
 Mathíldis Regínæ, matris Othónis Primi, Romanórum Imperatóris, humilitáte et 
 patiéntia conspícuæ.
\switchcolumn
\selectlanguage{english}
At Halberstadt in Germany, the death of blessed Queen Matilda, mother of 
 Emperor Otto I, renowned for humility and patience.
\switchcolumn*
\selectlanguage{latin}
\end{paracol}


% ---- martyrology/mart03/mart0315.htm
\needspace{10\baselineskip}
\begin{paracol}{2}
\selectlanguage{latin}
\begin{center}{\color{gregoriocolor} Idibus Mártii. 
 Luna\dots\ }\end{center}
\switchcolumn
\selectlanguage{english}
\begin{center}{\color{gregoriocolor} The 15\textsuperscript{th} Day of March. 
 The\dots\ Day of the Moon.}\end{center}
\end{paracol}

\noindent\begin{tabularx}{\linewidth}{*{19}{>{\centering\arraybackslash}X}}
 \textcolor{gregoriocolor}{a} & \textcolor{gregoriocolor}{b} & \textcolor{gregoriocolor}{c} & \textcolor{gregoriocolor}{d} & \textcolor{gregoriocolor}{e} & \textcolor{gregoriocolor}{f} & \textcolor{gregoriocolor}{g} & \textcolor{gregoriocolor}{h} & \textcolor{gregoriocolor}{i} & \textcolor{gregoriocolor}{k} & \textcolor{gregoriocolor}{l} & \textcolor{gregoriocolor}{m} & \textcolor{gregoriocolor}{n} & \textcolor{gregoriocolor}{p} & \textcolor{gregoriocolor}{q} & \textcolor{gregoriocolor}{r} & \textcolor{gregoriocolor}{s} & \textcolor{gregoriocolor}{t} & \textcolor{gregoriocolor}{u} \\
 16 & 17 & 18 & 19 & 20 & 21 & 22 & 23 & 24 & 25 & 26 & 27 & 28 & 29 & 30 & 1 & 2 & 3 & 4 \\
\end{tabularx}
\vspace{0.5\baselineskip}
\noindent\begin{tabularx}{\linewidth}{*{12}{>{\centering\arraybackslash}X}}
 \textcolor{gregoriocolor}{A} & \textcolor{gregoriocolor}{B} & \textcolor{gregoriocolor}{C} & \textcolor{gregoriocolor}{D} & \textcolor{gregoriocolor}{E} & F & \textcolor{gregoriocolor}{F} & \textcolor{gregoriocolor}{G} & \textcolor{gregoriocolor}{H} & \textcolor{gregoriocolor}{M} & \textcolor{gregoriocolor}{N} & \textcolor{gregoriocolor}{P} \\
 5 & 6 & 7 & 8 & 9 & 10 & 10 & 11 & 12 & 13 & 14 & 15 \\
\end{tabularx}

\begin{paracol}{2}
\selectlanguage{latin}
\lettrine[lines=2]{C}{æsaréæ,} in Cappadócia, pássio sancti Longíni 
 mílitis, qui Dómini latus láncea perforásse perhibétur.
\switchcolumn
\selectlanguage{english}
\lettrine[lines=2]{A}{t} Caesarea in Cappadocia, the martyrdom of St. Longinus, the soldier who is 
 said to have pierced our Lord's side with a lance.
\switchcolumn*
\selectlanguage{latin}
Eódem die natális sancti Aristobúli, 
 Apostolórum discípuli, qui, cursu prædicatiónis perácto, martyrium 
 consummávit.
\switchcolumn
\selectlanguage{english}
The same day, the birthday of St. Aristobulus, a disciple of the apostles, 
 who completed by martyrdom a life spent in preaching the Gospel.
\switchcolumn*
\selectlanguage{latin}
In Hellespónto sancti Menígni fullónis, qui sub Décio Imperatóre passus est.
\switchcolumn
\selectlanguage{english}
In the Hellespont, St. Menignus, a dyer, who suffered under Decius.
\switchcolumn*
\selectlanguage{latin}
In Ægypto sancti Nicándri Mártyris, qui, cum 
 sanctórum Mártyrum relíquias studióse perquíreret, et ipse méruit éffici 
 Martyr, sub Diocletiáno Imperatóre.
\switchcolumn
\selectlanguage{english}
In Egypt, St. Nicander, martyr, who sought diligently for the remains of the 
 holy martyrs, and thus merited to be made a martyr himself, under Emperor 
 Diocletian.
\switchcolumn*
\selectlanguage{latin}
Córdubæ, in Hispánia, sanctæ Leocrítiæ, 
 Vírginis et Mártyris; quæ ob Christi fidem, in persecutióne Arábica, 
 divérsis cruciátibus afflícta et gládio decolláta est.
\switchcolumn
\selectlanguage{english}
At Cordova in Spain, St. Leocritia, virgin and martyr. She suffered 
 various cruel tortures and was beheaded for the faith of Christ during the 
 Arabian persecution.
\switchcolumn*
\selectlanguage{latin}
Thessalonícæ sanctæ Matrónæ, quæ, cum esset 
 ancílla cujúsdam Judǽæ, et occúlte Christum cóleret, ac furtívis oratiónibus 
 quotídie Ecclésiam frequentáret, a dómina sua est deprehénsa et 
 multiplíciter afflícta, atque novíssime, robústis fústibus usque ad mortem 
 cæsa, in confessióne Christi, incorrúptum Deo spíritum réddidit.
\switchcolumn
\selectlanguage{english}
At Thessalonica, St. Matrona, servant of a Jewess, who, worshipping Christ 
 secretly, and stealing away daily to pray in the church, was detected by her 
 mistress and subjected to many trials. Being at last beaten to death 
 with large clubs, she gave up her pure soul to God in confessing Christ.
\switchcolumn*
\selectlanguage{latin}
Reáte sancti Probi Epíscopi, cui moriénti Juvenális et Eleuthérius Mártyres 
 adfuérunt.
\switchcolumn
\selectlanguage{english}
At Rieti, the bishop St. Probus, at whose death the martyrs Juvenal and 
 Eleutherius were present.
\switchcolumn*
\selectlanguage{latin}
Vindobónæ, in Austria, sancti Cleméntis-Maríæ 
 Hofbauer, Sacerdótis proféssi Congregatiónis a sanctíssimo Redemptóre 
 nuncupátæ, plúrimis in Dei glória et animárum salúte promovénda ac dilatánda 
 ipsa Congregatióne exantlátis labóribus insígnis; quem, virtútibus et 
 miráculis clarum, Pius Décimus, Póntifex Máximus, in Sanctórum cánonem 
 rétulit.
\switchcolumn
\selectlanguage{english}
At Vienna in Austria, St. Clement Mary Hofbauer, a priest of the 
 Congregation of the Most Holy Redeemer, renowned for his great devotion in 
 promoting the glory of God and the salvation of souls, and in extending that 
 order. He was canonized by Pope Pius X.
\switchcolumn*
\selectlanguage{latin}
Apud Cápuam sancti Speciósi Mónachi, cujus ánimam (ut scribit beátus 
 Gregórius Papa) germánus ejus deférri vidit in cælum.
\switchcolumn
\selectlanguage{english}
At Capua, the monk St. Speciosus, whose soul his brother saw being 
 taken into heaven, as is recorded by Pope St. Gregory.
\switchcolumn*
\selectlanguage{latin}
Lutétiæ Parisiórum sanctæ Ludovícæ de Marillac, 
 víduæ Le Gras, Societátis Puellárum a Caritáte una cum sancto Vincéntio a 
 Paulo Fundatrícis, egénis sublevándis addictíssimæ, quam Pius Papa Undécimus 
 Sanctárum fastis accénsuit.
\switchcolumn
\selectlanguage{english}
At Paris, the birthday of St. Louise de Marillac, a widow of Le Gras, 
 co-founder with St. Vincent de Paul of the Society of the Daughters of 
 Charity. Outstanding for her virtues and miracles, her name was 
 inscribed on the roll of the saints by Pope Pius XI.
\switchcolumn*
\selectlanguage{latin}
\end{paracol}


% ---- martyrology/mart03/mart0316.htm
\needspace{10\baselineskip}
\begin{paracol}{2}
\selectlanguage{latin}
\begin{center}{\color{gregoriocolor} Décimo séptimo Kaléndas Aprilis. 
 Luna\dots\ }\end{center}
\switchcolumn
\selectlanguage{english}
\begin{center}{\color{gregoriocolor} The 16\textsuperscript{th} Day of March. 
 The\dots\ Day of the Moon.}\end{center}
\end{paracol}

\noindent\begin{tabularx}{\linewidth}{*{19}{>{\centering\arraybackslash}X}}
 \textcolor{gregoriocolor}{a} & \textcolor{gregoriocolor}{b} & \textcolor{gregoriocolor}{c} & \textcolor{gregoriocolor}{d} & \textcolor{gregoriocolor}{e} & \textcolor{gregoriocolor}{f} & \textcolor{gregoriocolor}{g} & \textcolor{gregoriocolor}{h} & \textcolor{gregoriocolor}{i} & \textcolor{gregoriocolor}{k} & \textcolor{gregoriocolor}{l} & \textcolor{gregoriocolor}{m} & \textcolor{gregoriocolor}{n} & \textcolor{gregoriocolor}{p} & \textcolor{gregoriocolor}{q} & \textcolor{gregoriocolor}{r} & \textcolor{gregoriocolor}{s} & \textcolor{gregoriocolor}{t} & \textcolor{gregoriocolor}{u} \\
 17 & 18 & 19 & 20 & 21 & 22 & 23 & 24 & 25 & 26 & 27 & 28 & 29 & 30 & 1 & 2 & 3 & 4 & 5 \\
\end{tabularx}
\vspace{0.5\baselineskip}
\noindent\begin{tabularx}{\linewidth}{*{12}{>{\centering\arraybackslash}X}}
 \textcolor{gregoriocolor}{A} & \textcolor{gregoriocolor}{B} & \textcolor{gregoriocolor}{C} & \textcolor{gregoriocolor}{D} & \textcolor{gregoriocolor}{E} & F & \textcolor{gregoriocolor}{F} & \textcolor{gregoriocolor}{G} & \textcolor{gregoriocolor}{H} & \textcolor{gregoriocolor}{M} & \textcolor{gregoriocolor}{N} & \textcolor{gregoriocolor}{P} \\
 6 & 7 & 8 & 9 & 10 & 11 & 11 & 12 & 13 & 14 & 15 & 16 \\
\end{tabularx}

\begin{paracol}{2}
\selectlanguage{latin}
\lettrine[lines=2]{R}{omæ} pássio sancti Cyríaci Diáconi, qui, post 
 longam cárceris maceratiónem, liquáta pice perfúsus et in catásta exténsus, 
 attráctus étiam nervis et fústibus cæsus, ad últimum, cum Largo et Smarágdo 
 et áliis vigínti, jubénte Maximiáno, cápite truncátus est. Sanctórum 
 vero Cyríaci, Largi et Smarágdi festívitas sexto Idus Augústi recólitur, quo 
 die a beáto Marcéllo Papa córpora eorúndem vigínti trium Mártyrum leváta 
 sunt ac venerabíliter tumuláta.
\switchcolumn
\selectlanguage{english}
\lettrine[lines=2]{A}{t} Rome the martyrdom of the deacon St. Cyriacus, who, after a long 
 imprisonment, had melted pitch poured over him, was stretched on the rack, 
 had his limbs pulled with ropes, was beaten with clubs, and finally was 
 beheaded by order of Maximian, together with Largus, Smaragdus, and twenty 
 others. Their feast, however, is kept on the 8th of August, the day on 
 which these twenty-three martyrs were exhumed by blessed Pope Marcellus and 
 reverently entombed.
\switchcolumn*
\selectlanguage{latin}
Aquiléjæ natális beáti Hilárii Epíscopi, et 
 Tatiáni Diáconi, qui, sub Numeriáno Imperatóre et Berónio Præside, post 
 equúleum atque ália torménta, una cum Felíce, Largo et Dionysio, martyrium 
 terminárunt.
\switchcolumn
\selectlanguage{english}
At Aquileia, in the time of Emperor Numerian and the governor Beronius, the 
 birthday of the holy bishop Hilary, and the deacon Tatian, who were martyred 
 with Felix, Largus, and Denis, after being subjected to the rack and other 
 tortures.
\switchcolumn*
\selectlanguage{latin}
In Lycaónia sancti Papæ Mártyris, qui, ob 
 Christi fidem, verbéribus cæsus, úngulis férreis lacerátus, clavátis cálceis 
 incédere jussus est; deínde árbori alligátus, eándem árborem, migrans ad 
 Dóminum, ex stérili réddidit fructuósam.
\switchcolumn
\selectlanguage{english}
In Lycaonia, the holy martyr Papas, who was scourged for the Christian 
 faith, had his flesh torn with iron hooks, and was compelled to walk with 
 shoes pierced with nails, and was finally bound to a barren tree. In 
 leaving this world to go to God, he rendered this same tree fruitful.
\switchcolumn*
\selectlanguage{latin}
Anazárbi, in Cilícia, sancti Juliáni Mártyris, qui, sub Marciáno Præside, 
 diutíssime cruciátus, demum, in sacco una cum serpéntibus inclúsus, in mare 
 demérsus est.
\switchcolumn
\selectlanguage{english}
At Anazarbum in Cilicia, under the governor Marcian, the martyr St. Julian, 
 who was a long time tortured, then put into a sack with serpents, and cast 
 into the sea.
\switchcolumn*
\selectlanguage{latin}
In dicióne Canadénsi sanctórum Mártyrum Joánnis de Brébeuf, 
 Gabriélis Lalemant, Antónii Daniel, Cároli Garnier et Natális Chabanel, 
 Presbyterórum Societátis Jesu, qui in Hurónica Missióne, hac aliísque diébus, 
 post multos labóres et sævíssimos cruciátus, mortem pro Christo fórtiter 
 obiérunt.
\switchcolumn
\selectlanguage{english}
In the territory of Canada, Saints John de Brébeuf, Gabriel Lalemant, 
 Anthony Daniel, Charles Garnier, and Noel Chabanel, priests of the Society 
 of Jesus, who in the mission of the Hurons, on this and other days, after 
 many labours and most cruel torments, bravely underwent death for Christ.
\switchcolumn*
\selectlanguage{latin}
Ravénnæ sancti Agapíti, Epíscopi et Confessóris.
\switchcolumn
\selectlanguage{english}
At Ravenna, St. Agapitus, bishop and confessor.
\switchcolumn*
\selectlanguage{latin}
Colóniæ Agrippínæ sancti Heribérti Epíscopi, 
 sanctitáte célebris.
\switchcolumn
\selectlanguage{english}
At Cologne, St. Heribert, bishop, celebrated for sanctity.
\switchcolumn*
\selectlanguage{latin}
Arvérnis, in Gállia, deposítio sancti Patrícii Epíscopi.
\switchcolumn
\selectlanguage{english}
In Auvergne, the death of St. Patrick, bishop.
\switchcolumn*
\selectlanguage{latin}
In Syria sancti Abrahæ Eremítæ, cujus res 
 gestas beátus Ephræm Diáconus conscrípsit.
\switchcolumn
\selectlanguage{english}
In Syria, St. Abraham, hermit, whose life has been written by the blessed 
 deacon Ephrem.
\switchcolumn*
\selectlanguage{latin}
\end{paracol}


% ---- martyrology/mart03/mart0317.htm
\needspace{10\baselineskip}
\begin{paracol}{2}
\selectlanguage{latin}
\begin{center}{\color{gregoriocolor} Sextodécimo Kaléndas Aprilis. 
 Luna\dots\ }\end{center}
\switchcolumn
\selectlanguage{english}
\begin{center}{\color{gregoriocolor} The 17\textsuperscript{th} Day of March. 
 The\dots\ Day of the Moon.}\end{center}
\end{paracol}

\noindent\begin{tabularx}{\linewidth}{*{19}{>{\centering\arraybackslash}X}}
 \textcolor{gregoriocolor}{a} & \textcolor{gregoriocolor}{b} & \textcolor{gregoriocolor}{c} & \textcolor{gregoriocolor}{d} & \textcolor{gregoriocolor}{e} & \textcolor{gregoriocolor}{f} & \textcolor{gregoriocolor}{g} & \textcolor{gregoriocolor}{h} & \textcolor{gregoriocolor}{i} & \textcolor{gregoriocolor}{k} & \textcolor{gregoriocolor}{l} & \textcolor{gregoriocolor}{m} & \textcolor{gregoriocolor}{n} & \textcolor{gregoriocolor}{p} & \textcolor{gregoriocolor}{q} & \textcolor{gregoriocolor}{r} & \textcolor{gregoriocolor}{s} & \textcolor{gregoriocolor}{t} & \textcolor{gregoriocolor}{u} \\
 18 & 19 & 20 & 21 & 22 & 23 & 24 & 25 & 26 & 27 & 28 & 29 & 30 & 1 & 2 & 3 & 4 & 5 & 6 \\
\end{tabularx}
\vspace{0.5\baselineskip}
\noindent\begin{tabularx}{\linewidth}{*{12}{>{\centering\arraybackslash}X}}
 \textcolor{gregoriocolor}{A} & \textcolor{gregoriocolor}{B} & \textcolor{gregoriocolor}{C} & \textcolor{gregoriocolor}{D} & \textcolor{gregoriocolor}{E} & F & \textcolor{gregoriocolor}{F} & \textcolor{gregoriocolor}{G} & \textcolor{gregoriocolor}{H} & \textcolor{gregoriocolor}{M} & \textcolor{gregoriocolor}{N} & \textcolor{gregoriocolor}{P} \\
 7 & 8 & 9 & 10 & 11 & 12 & 12 & 13 & 14 & 15 & 16 & 17 \\
\end{tabularx}

\begin{paracol}{2}
\selectlanguage{latin}
\lettrine[lines=2]{A}{pud} civitátem Dunum, in Hibérnia, natális sancti Patrícii, Epíscopi et 
 Confessóris, qui primus in ea ínsula Christum evangelizávit, et máximis 
 miráculis et virtútibus cláruit.
\switchcolumn
\selectlanguage{english}
\lettrine[lines=2]{A}{t} Downpatrick in Ireland, the birthday of St. Patrick, bishop and 
 confessor, who was the first to preach Christ in that country, and who 
 became illustrious by great miracles and virtues.
\switchcolumn*
\selectlanguage{latin}
Hierosólymis sancti Joseph ab Arimathǽa, qui 
 nóbilis Decúrio et discípulus Dómini éxstitit; atque ipsíus Dómini corpus, 
 de cruce depósitum, in monuménto suo novo sepelívit.
\switchcolumn
\selectlanguage{english}
At Jerusalem, St. Joseph of Arimathea, noble senator and disciple of our 
 Lord,. who took his Body down from the Cross and buried it in his own new 
 sepulchre.
\switchcolumn*
\selectlanguage{latin}
Romæ sanctórum Alexándri et Theodóri Mártyrum.
\switchcolumn
\selectlanguage{english}
At Rome, the Saints Alexander and Theodore, martyrs.
\switchcolumn*
\selectlanguage{latin}
Alexandríæ commemorátio plurimórum sanctórum 
 Mártyrum, qui a Serápidis cultóribus comprehénsi, et, cum adoráre idólum 
 constánter renuíssent, sævíssime cæsi sunt, témpore Theodósii Imperatóris; 
 qui mox rescríptum dedit, ut Serápidis templum destruerétur.
\switchcolumn
\selectlanguage{english}
At Alexandria, the commemoration of many holy martyrs, who, being seized by 
 the worshippers of Serapis, and refusing constantly to adore that idol, were 
 cruelly murdered. Emperor Theodosius, who issued the order, afterwards 
 commanded that the temple of Serapis should be destroyed.
\switchcolumn*
\selectlanguage{latin}
Constantinópoli sancti Pauli Mártyris, qui, sub Constantíno Coprónymo, 
 cum sanctárum Imáginum cultum defénderet, igne combústus est.
\switchcolumn
\selectlanguage{english}
At Constantinople, St. Paul, martyr, who was burned alive by Constantine 
 Copronymus, for defending the veneration of sacred images.
\switchcolumn*
\selectlanguage{latin}
Cabillóne, in Gálliis, sancti Agrícolæ Epíscopi.
\switchcolumn
\selectlanguage{english}
At Chalons in France, St. Agricola, bishop.
\switchcolumn*
\selectlanguage{latin}
Nivigéllæ, in Brabántia, sanctæ Gertrúdis 
 Vírginis, quæ claríssimo génere orta, despíciens mundum et toto vitæ suæ 
 cursu in ómnibus sanctitátis offíciis se exércens, Christum sponsum in cælis 
 habére méruit.
\switchcolumn
\selectlanguage{english}
At Nivelle in Brabant, St. Gertrude, a virgin of noble birth. Because 
 she despised the world, and during her whole life practised all kinds of 
 good works, she deserved to have Christ for her spouse in heaven.
\switchcolumn*
\selectlanguage{latin}
\end{paracol}


% ---- martyrology/mart03/mart0318.htm
\needspace{10\baselineskip}
\begin{paracol}{2}
\selectlanguage{latin}
\begin{center}{\color{gregoriocolor} Quintodécimo Kaléndas Aprilis. 
 Luna\dots\ }\end{center}
\switchcolumn
\selectlanguage{english}
\begin{center}{\color{gregoriocolor} The 18\textsuperscript{th} Day of March. 
 The\dots\ Day of the Moon.}\end{center}
\end{paracol}

\noindent\begin{tabularx}{\linewidth}{*{19}{>{\centering\arraybackslash}X}}
 \textcolor{gregoriocolor}{a} & \textcolor{gregoriocolor}{b} & \textcolor{gregoriocolor}{c} & \textcolor{gregoriocolor}{d} & \textcolor{gregoriocolor}{e} & \textcolor{gregoriocolor}{f} & \textcolor{gregoriocolor}{g} & \textcolor{gregoriocolor}{h} & \textcolor{gregoriocolor}{i} & \textcolor{gregoriocolor}{k} & \textcolor{gregoriocolor}{l} & \textcolor{gregoriocolor}{m} & \textcolor{gregoriocolor}{n} & \textcolor{gregoriocolor}{p} & \textcolor{gregoriocolor}{q} & \textcolor{gregoriocolor}{r} & \textcolor{gregoriocolor}{s} & \textcolor{gregoriocolor}{t} & \textcolor{gregoriocolor}{u} \\
 19 & 20 & 21 & 22 & 23 & 24 & 25 & 26 & 27 & 28 & 29 & 30 & 1 & 2 & 3 & 4 & 5 & 6 & 7 \\
\end{tabularx}
\vspace{0.5\baselineskip}
\noindent\begin{tabularx}{\linewidth}{*{12}{>{\centering\arraybackslash}X}}
 \textcolor{gregoriocolor}{A} & \textcolor{gregoriocolor}{B} & \textcolor{gregoriocolor}{C} & \textcolor{gregoriocolor}{D} & \textcolor{gregoriocolor}{E} & F & \textcolor{gregoriocolor}{F} & \textcolor{gregoriocolor}{G} & \textcolor{gregoriocolor}{H} & \textcolor{gregoriocolor}{M} & \textcolor{gregoriocolor}{N} & \textcolor{gregoriocolor}{P} \\
 8 & 9 & 10 & 11 & 12 & 13 & 13 & 14 & 15 & 16 & 17 & 18 \\
\end{tabularx}

\begin{paracol}{2}
\selectlanguage{latin}
\lettrine[lines=2]{H}{ierosólymis} sancti Cyrílli Epíscopi, Confessóris et Ecclésiæ 
 Doctóris; qui, ab Ariánis multas pro fidei causa perpéssus injúrias et ex 
 Ecclésia sua sæpe depúlsus, tandem, sanctitátis glória clarus, in pace 
 quiévit. Ipsíus porro intemerátam fidem prima Constantinopolitána 
 Synodus œcuménica, sancto Dámaso Papæ scribens, præcláro testimónio 
 commendávit.
\switchcolumn
\selectlanguage{english}
\lettrine[lines=2]{A}{t} Jerusalem, St. Cyril, bishop, who suffered many injuries from the Arians 
 for the faith. Often exiled from his church, he at length rested in 
 peace with a great reputation for sanctity. A magnificent testimony of 
 the purity of his faith is given by the first ecumenical Council of 
 Constantinople in a letter to Pope Damasus.
\switchcolumn*
\selectlanguage{latin}
Cæsaréæ, in Palæstína, natális beáti Alexándri 
 Epíscopi, qui de Cappadócia, ex própria civitáte, ubi erat Epíscopus, 
 sanctórum locórum desidério Hierosólymam pétiit; atque ibi, cum a Narcísso, 
 ejúsdem urbis Epíscopo, jam sene, illa regerétur Ecclésia, ipsíus 
 gubernácula, divína edóctus revelatióne, suscépit. Póstmodum vero, in 
 persecutióne Décii, cum jam longævæ ætátis veneránda canítie præfúlgeret, 
 ductus est Cæsaréam, et clausus in cárcere, ob confessiónem Christi, 
 martyrium complévit.
\switchcolumn
\selectlanguage{english}
At Caesarea in Palestine, the birthday of the blessed Bishop Alexander, who, 
 from his own city in Cappadocia, where he was bishop, coming to Jerusalem to 
 visit the holy places, took upon himself, by divine revelation, the 
 government of that church in place of the aged Narcissus. Sometime 
 afterwards, when he had become venerable by his age and gray hair, he was 
 led to Caesarea and shut up in prison, where he completed his martyrdom for 
 the confession of Christ during the persecution of Decius.
\switchcolumn*
\selectlanguage{latin}
Augústæ sancti Narcíssi Epíscopi, qui primus in 
 Rhǽtia Evangélium prædicávit; deínde in Hispániam profectus est, et, cum 
 Gerúndæ multos ad Christi fidem convertísset, ibídem, in persecutióne 
 Diocletiáni Imperatóris, una cum Felíce Diácono, martyrii palmam accépit.
\switchcolumn
\selectlanguage{english}
At Augsburg, St. Narcissus, bishop, who was the first to preach the Gospel 
 in the Tyrol. Afterwards, setting out for Spain, he converted many to 
 the faith of Christ at Gerona, and there, along with the deacon Felix, he 
 received the palm of martyrdom during the persecution of Diocletian.
\switchcolumn*
\selectlanguage{latin}
Nicomedíæ sanctórum decem míllium Mártyrum, 
 qui, pro Christi confessióne, gládio percússi sunt.
\switchcolumn
\selectlanguage{english}
At Nicomedia, ten thousand holy martyrs, who were put to the sword for the 
 confession of Christ.
\switchcolumn*
\selectlanguage{latin}
Ibídem sanctórum Mártyrum Tróphimi et Eucárpii.
\switchcolumn
\selectlanguage{english}
In the same place, the holy martyrs Trophimus and Eucarpius.
\switchcolumn*
\selectlanguage{latin}
In Británnia sancti Eduárdi Regis, qui, dolis novércæ 
 necátus, multis miráculis cláruit.
\switchcolumn
\selectlanguage{english}
In England, St. Edward, king, who was assassinated by order of his 
 treacherous stepmother, and became celebrated for many miracles.
\switchcolumn*
\selectlanguage{latin}
Lucæ, in Túscia, natális sancti Frigdiáni 
 Epíscopi, virtúte miraculórum illústris.
\switchcolumn
\selectlanguage{english}
At Lucca in Tuscany, the birthday of the holy bishop Fridian, who was 
 illustrious by the power of working miracles.
\switchcolumn*
\selectlanguage{latin}
Mántuæ sancti Ansélmi, Epíscopi Lucénsis et 
 Confessóris.
\switchcolumn
\selectlanguage{english}
At Mantua, St. Anselm, bishop and confessor.
\switchcolumn*
\selectlanguage{latin}
Cárali, in Sardínia, sancti Salvatóris ab Horta 
 Confessóris, ex Ordine Fratrum Minórum, qui virtútibus et singulári 
 miraculórum dono cláruit, et a Pio Papa Undécimo inter sanctos Cǽlites 
 adnumerátus est.
\switchcolumn
\selectlanguage{english}
At Cagliari in Sardinia, St. Salvatore of Orte, confessor, a member of the 
 Order of Friars Minor, who was numbered among the heavenly saints by Pope 
 Pius XI, because he was graced with every virtue and had been given by God 
 the gift of performing outstanding miracles.
\switchcolumn*
\selectlanguage{latin}
\end{paracol}


% ---- martyrology/mart03/mart0319.htm
\needspace{10\baselineskip}
\begin{paracol}{2}
\selectlanguage{latin}
\begin{center}{\color{gregoriocolor} Quartodécimo Kaléndas Aprilis. 
 Luna\dots\ }\end{center}
\switchcolumn
\selectlanguage{english}
\begin{center}{\color{gregoriocolor} The 19\textsuperscript{th} Day of March. 
 The\dots\ Day of the Moon.}\end{center}
\end{paracol}

\noindent\begin{tabularx}{\linewidth}{*{19}{>{\centering\arraybackslash}X}}
 \textcolor{gregoriocolor}{a} & \textcolor{gregoriocolor}{b} & \textcolor{gregoriocolor}{c} & \textcolor{gregoriocolor}{d} & \textcolor{gregoriocolor}{e} & \textcolor{gregoriocolor}{f} & \textcolor{gregoriocolor}{g} & \textcolor{gregoriocolor}{h} & \textcolor{gregoriocolor}{i} & \textcolor{gregoriocolor}{k} & \textcolor{gregoriocolor}{l} & \textcolor{gregoriocolor}{m} & \textcolor{gregoriocolor}{n} & \textcolor{gregoriocolor}{p} & \textcolor{gregoriocolor}{q} & \textcolor{gregoriocolor}{r} & \textcolor{gregoriocolor}{s} & \textcolor{gregoriocolor}{t} & \textcolor{gregoriocolor}{u} \\
 20 & 21 & 22 & 23 & 24 & 25 & 26 & 27 & 28 & 29 & 30 & 1 & 2 & 3 & 4 & 5 & 6 & 7 & 8 \\
\end{tabularx}
\vspace{0.5\baselineskip}
\noindent\begin{tabularx}{\linewidth}{*{12}{>{\centering\arraybackslash}X}}
 \textcolor{gregoriocolor}{A} & \textcolor{gregoriocolor}{B} & \textcolor{gregoriocolor}{C} & \textcolor{gregoriocolor}{D} & \textcolor{gregoriocolor}{E} & F & \textcolor{gregoriocolor}{F} & \textcolor{gregoriocolor}{G} & \textcolor{gregoriocolor}{H} & \textcolor{gregoriocolor}{M} & \textcolor{gregoriocolor}{N} & \textcolor{gregoriocolor}{P} \\
 9 & 10 & 11 & 12 & 13 & 14 & 14 & 15 & 16 & 17 & 18 & 19 \\
\end{tabularx}

\begin{paracol}{2}
\selectlanguage{latin}
\lettrine[lines=2]{I}{n} Judæa natális sancti Joseph, Sponsi 
 beatíssimæ Vírginis Maríæ, Confessóris; quem Pius Nonus, Póntifex Máximus, 
 votis et précibus ánnuens totíus cathólici Orbis, universális Ecclésiæ 
 Patrónum declarávit.
\switchcolumn
\selectlanguage{english}
\lettrine[lines=2]{I}{n} Judea, the birthday of St. Joseph, spouse of the Most Blessed Virgin 
 Mary, Confessor and Patron of the Universal Church. Pope Pius IX, yielding to the desires and prayers of the whole 
 Catholic world, declared him Patron of the Universal Church.
\switchcolumn*
\selectlanguage{latin}
Surrénti sanctórum Mártyrum Quincti, Quinctíllæ, 
 Quartíllæ et Marci, cum áliis novem.
\switchcolumn
\selectlanguage{english}
At Sorrento, the holy martyrs Quinctus, Quinctilla, Quartilla, Mark, and 
 nine others.
\switchcolumn*
\selectlanguage{latin}
Nicomedíæ sancti Panchárii Románi, qui, sub 
 Diocletiáno Imperatóre, in hujus grátiam Christum pro diis inánibus ejurávit, 
 sed, matre ac soróre instántibus, ad veram fidem mox redívit, et ob immótam 
 in ea constántiam, nervis cæsus, et cápite truncátus, martyrii corónam 
 accépit.
\switchcolumn
\selectlanguage{english}
At Nicomedia, St. Pancharius, a Roman, who apostatized for the sake of 
 Emperor Diocletian, but by the persuasion of his mother and sister 
 immediately returned to the true faith. Because of his subsequent 
 constancy in it, he was beaten with clubs and beheaded, obtaining thus the 
 crown of martyrdom.
\switchcolumn*
\selectlanguage{latin}
Eódem die sanctórum Apollónii et Leóntii Episcopórum.
\switchcolumn
\selectlanguage{english}
The same day, the holy Bishops Apollonius and Leontius.
\switchcolumn*
\selectlanguage{latin}
Gandávi, in Flándria, sanctórum Landoáldi, Presbyteri Románi, et Amántii 
 Diáconi; qui, a sancto Martíno Papa ad prædicándum 
 Evangélium missi, ambo apostólicum sibi commíssum opus fidéliter implevérunt, 
 ac multis post óbitum sunt illustráti miráculis.
\switchcolumn
\selectlanguage{english}
At Ghent in Flanders, Saints Landoald, a Roman priest, and the deacon 
 Amantius, who were sent to preach the Gospel by Pope St. Martin. They 
 faithfully fulfilled this apostolic appointment, and after their deaths 
 became renowned for their miracles.
\switchcolumn*
\selectlanguage{latin}
Apud Pinnénsem civitátem natális beáti Joánnis, magnæ 
 sanctitátis viri; qui de Syria ad Itáliam venit, atque ibi, constrúcto 
 monastério, multórum servórum Dei per quátuor et quadragínta annos Pater éxstitit, et, clarus virtútibus, in pace quiévit.
\switchcolumn
\selectlanguage{english}
In the city of Pinna, the birthday of blessed John, a man of great sanctity, 
 who came from Syria into Italy, and there founded a monastery. After 
 being the spiritual guide for many of God's servants for forty-four years, 
 he rested in peace.
\switchcolumn*
\selectlanguage{latin}
\end{paracol}


% ---- martyrology/mart03/mart0320.htm
\needspace{10\baselineskip}
\begin{paracol}{2}
\selectlanguage{latin}
\begin{center}{\color{gregoriocolor} Tertiodécimo Kaléndas Aprilis. 
 Luna\dots\ }\end{center}
\switchcolumn
\selectlanguage{english}
\begin{center}{\color{gregoriocolor} The 20\textsuperscript{th} Day of March. 
 The\dots\ Day of the Moon.}\end{center}
\end{paracol}

\noindent\begin{tabularx}{\linewidth}{*{19}{>{\centering\arraybackslash}X}}
 \textcolor{gregoriocolor}{a} & \textcolor{gregoriocolor}{b} & \textcolor{gregoriocolor}{c} & \textcolor{gregoriocolor}{d} & \textcolor{gregoriocolor}{e} & \textcolor{gregoriocolor}{f} & \textcolor{gregoriocolor}{g} & \textcolor{gregoriocolor}{h} & \textcolor{gregoriocolor}{i} & \textcolor{gregoriocolor}{k} & \textcolor{gregoriocolor}{l} & \textcolor{gregoriocolor}{m} & \textcolor{gregoriocolor}{n} & \textcolor{gregoriocolor}{p} & \textcolor{gregoriocolor}{q} & \textcolor{gregoriocolor}{r} & \textcolor{gregoriocolor}{s} & \textcolor{gregoriocolor}{t} & \textcolor{gregoriocolor}{u} \\
 21 & 22 & 23 & 24 & 25 & 26 & 27 & 28 & 29 & 30 & 1 & 2 & 3 & 4 & 5 & 6 & 7 & 8 & 9 \\
\end{tabularx}
\vspace{0.5\baselineskip}
\noindent\begin{tabularx}{\linewidth}{*{12}{>{\centering\arraybackslash}X}}
 \textcolor{gregoriocolor}{A} & \textcolor{gregoriocolor}{B} & \textcolor{gregoriocolor}{C} & \textcolor{gregoriocolor}{D} & \textcolor{gregoriocolor}{E} & F & \textcolor{gregoriocolor}{F} & \textcolor{gregoriocolor}{G} & \textcolor{gregoriocolor}{H} & \textcolor{gregoriocolor}{M} & \textcolor{gregoriocolor}{N} & \textcolor{gregoriocolor}{P} \\
 10 & 11 & 12 & 13 & 14 & 15 & 15 & 16 & 17 & 18 & 19 & 20 \\
\end{tabularx}

\begin{paracol}{2}
\selectlanguage{latin}
\lettrine[lines=2]{I}{n} Judǽa natális sancti Jóachim, patris 
 immaculátæ Vírginis Genitrícis Dei Maríæ, Confessóris. Ipsíus tamen 
 festum agitur décimo séptimo Kaléndas Septémbris.
\switchcolumn
\selectlanguage{english}
\lettrine[lines=2]{I}{n} Judea, St. Joachim, the father of the Most Blessed Virgin Mary, Mother of 
 God. His feast day is on the 16th of August.
\switchcolumn*
\selectlanguage{latin}
In Asia item natális sancti Archíppi, qui beáti Pauli Apóstoli 
 éxstitit 
 commílito, et cujus ipse in Epístola ad Philémonem 
 et ad Colossénses méminit.
\switchcolumn
\selectlanguage{english}
In Asia, the birthday of St. Archippus, fellow-labourer of the apostle St. 
 Paul, who is mentioned by him in his epistles to Philemon and the 
 Colossians.
\switchcolumn*
\selectlanguage{latin}
In Syria sanctórum Mártyrum Pauli, Cyrílli, Eugénii et aliórum quátuor.
\switchcolumn
\selectlanguage{english}
In Syria, the holy martyrs Paul, Cyril, Eugene, and four others.
\switchcolumn*
\selectlanguage{latin}
Eódem die sanctórum Photínæ Samaritánæ, Joseph 
 et Victóris filiórum, itémque Sebastiáni Ducis, Anatólii, Phótii, Phótidis, 
 Parascéves et Cyríacæ germanárum; qui omnes, Christum conféssi, martyrium 
 sunt assecúti.
\switchcolumn
\selectlanguage{english}
On the same day, the Saints Photina, a Samaritan, and her sons Joseph and 
 Victor; also, Sebastian, a military officer, Anatolius, and Photius; 
 Photides, Parasceves, and Cyriaca, sisters, all of whom were put to death 
 for the confession of the faith.
\switchcolumn*
\selectlanguage{latin}
Amísi, in Paphlagónia, sanctárum septem mulíerum, scílicet Alexándræ, 
 Cláudiæ, Euphrásiæ, Matrónæ, Juliánæ, Euphémiæ et Theodósiæ; quæ in fidei 
 confessióne sunt cæsæ, eásque secútæ sunt Derphúta et soror ipsíus.
\switchcolumn
\selectlanguage{english}
At Amisus in Paphlagonia, seven holy women, Alexandra, Claudia, Euphrasia, 
 Matrona, Juliana, Euphemia, and Theodosia, who were put to death for the 
 confession of the faith. They were followed by Dephuta and her sister.
\switchcolumn*
\selectlanguage{latin}
Apollóniæ sancti Nicétæ Epíscopi, qui, pro 
 sanctárum Imáginum cultu ejéctus in exsílium, illic réddidit spíritum.
\switchcolumn
\selectlanguage{english}
At Apollonia, Bishop St. Nicetas, who died in exile where he had been sent 
 for upholding the veneration of sacred images.
\switchcolumn*
\selectlanguage{latin}
In monastério Fontanéllæ, in Gállia, sancti 
 Wulfránni, Epíscopi Senonénsis, qui, relícto Episcopátu, ibídem, clarus 
 miráculis, decéssit e vita.
\switchcolumn
\selectlanguage{english}
In the monastery of Fontanelle in France, St. Wulfram, bishop of Sens, who 
 resigned his bishopric, and after having performed miracles, departed out of 
 this life.
\switchcolumn*
\selectlanguage{latin}
In Británnia deposítio sancti Cuthbérti, Epíscopi Lindisfarnénsis, qui, a 
 puerítia ad óbitum usque, sanctis opéribus et miraculórum signis effúlsit.
\switchcolumn
\selectlanguage{english}
In England, the death of St. Cuthbert, bishop of Lindisfarne, who from his 
 childhood until his death was renowned for good works and miracles.
\switchcolumn*
\selectlanguage{latin}
Senis, in Túscia, Beáti Ambrósii, ex Ordine Prædicatórum, 
 sanctitáte, prædicatióne et miráculis clari.
\switchcolumn
\selectlanguage{english}
At Sienna in Tuscany, blessed Ambrose of the Order of Preachers, celebrated 
 for sanctity, eloquence, and miracles.
\switchcolumn*
\selectlanguage{latin}
\end{paracol}


% ---- martyrology/mart03/mart0321.htm
\needspace{10\baselineskip}
\begin{paracol}{2}
\selectlanguage{latin}
\begin{center}{\color{gregoriocolor} Duodécimo Kaléndas Aprílis. 
 Luna\dots\ }\end{center}
\switchcolumn
\selectlanguage{english}
\begin{center}{\color{gregoriocolor} The 21\textsuperscript{st} Day of March. 
 The\dots\ Day of the Moon.}\end{center}
\end{paracol}

\noindent\begin{tabularx}{\linewidth}{*{19}{>{\centering\arraybackslash}X}}
 \textcolor{gregoriocolor}{a} & \textcolor{gregoriocolor}{b} & \textcolor{gregoriocolor}{c} & \textcolor{gregoriocolor}{d} & \textcolor{gregoriocolor}{e} & \textcolor{gregoriocolor}{f} & \textcolor{gregoriocolor}{g} & \textcolor{gregoriocolor}{h} & \textcolor{gregoriocolor}{i} & \textcolor{gregoriocolor}{k} & \textcolor{gregoriocolor}{l} & \textcolor{gregoriocolor}{m} & \textcolor{gregoriocolor}{n} & \textcolor{gregoriocolor}{p} & \textcolor{gregoriocolor}{q} & \textcolor{gregoriocolor}{r} & \textcolor{gregoriocolor}{s} & \textcolor{gregoriocolor}{t} & \textcolor{gregoriocolor}{u} \\
 22 & 23 & 24 & 25 & 26 & 27 & 28 & 29 & 30 & 1 & 2 & 3 & 4 & 5 & 6 & 7 & 8 & 9 & 10 \\
\end{tabularx}
\vspace{0.5\baselineskip}
\noindent\begin{tabularx}{\linewidth}{*{12}{>{\centering\arraybackslash}X}}
 \textcolor{gregoriocolor}{A} & \textcolor{gregoriocolor}{B} & \textcolor{gregoriocolor}{C} & \textcolor{gregoriocolor}{D} & \textcolor{gregoriocolor}{E} & F & \textcolor{gregoriocolor}{F} & \textcolor{gregoriocolor}{G} & \textcolor{gregoriocolor}{H} & \textcolor{gregoriocolor}{M} & \textcolor{gregoriocolor}{N} & \textcolor{gregoriocolor}{P} \\
 11 & 12 & 13 & 14 & 15 & 16 & 16 & 17 & 18 & 19 & 20 & 21 \\
\end{tabularx}

\begin{paracol}{2}
\selectlanguage{latin}
\lettrine[lines=2]{I}{n} monte Cassíno natális sancti Benedícti Abbátis, qui in Occidénte fere 
 collápsam Monachórum disciplínam restítuit ac mirífice propagávit; cujus 
 vitam, virtútibus et miráculis gloriósam, beátus Gregórius Papa conscrípsit.
\switchcolumn
\selectlanguage{english}
\lettrine[lines=2]{A}{t} Monte Cassino, the birthday of the holy abbot St. Benedict, who restored 
 and wonderfully extended the monastic discipline in the West, where it had 
 almost been destroyed. His life, brilliant in virtues and miracles, 
 was written by Pope St. Gregory.
\switchcolumn*
\selectlanguage{latin}
Cátanæ, in Sicília, sancti Birílli, qui, a 
 beáto Petro ordinátus Epíscopus, ibídem, cum multos Gentílium convertísset 
 ad fidem, in última senectúte quiévit in pace.
\switchcolumn
\selectlanguage{english}
At Catania, St. Birillus, who was consecrated bishop by St. Peter. 
 After converting many gentiles to the faith, he died in extreme old age.
\switchcolumn*
\selectlanguage{latin}
Alexandríæ commemorátio sanctórum Mártyrum, 
 qui, sub Constántio Imperatóre et Præfécto Philágrio, irruéntibus Ariánis et 
 Gentílibus in Ecclésias, in die Parascéves cæsi sunt.
\switchcolumn
\selectlanguage{english}
At Alexandria, under Emperor Constantine and the governor Philagrius, the 
 commemoration of the holy martyrs who were murdered by the Arians and the 
 heathens, being attacked by them while they were in church on Good Friday.
\switchcolumn*
\selectlanguage{latin}
Eódem die sanctórum Mártyrum Philémonis et 
 Domníni.
\switchcolumn
\selectlanguage{english}
On the same day, the holy martyrs Philemon and Domninus.
\switchcolumn*
\selectlanguage{latin}
Alexandríæ beáti Serapiónis, Anachorétæ et 
 Epíscopi Thmúeos, magnárum virtútum viri; qui, Arianórum furóre in exsílium 
 actus, Conféssor migrávit ad Dóminum.
\switchcolumn
\selectlanguage{english}
At Alexandria, blessed Serapion, anchoret and bishop of Thmuis, a man of 
 great virtue, who was driven into exile by the enraged Arians, where he 
 departed to heaven.
\switchcolumn*
\selectlanguage{latin}
In território Lugdunénsi sancti Lupicíni Abbátis, cujus vita ob sanctitátis 
 et miraculórum glóriam fuit illústris.
\switchcolumn
\selectlanguage{english}
In the territory of Lyons, St. Lupicinus, abbot, whose life was brilliant 
 with the glory of holiness and miracles.
\switchcolumn*
\selectlanguage{latin}
In loco Ranft, prope Sachseln, in Helvétia, sancti Nicolái de Flüe, 
 patris famílias, dein Anachorétæ, arctíssima pæniténtia et mundi contémptu 
 insígnis, ab Helvétiis pater pátriæ appelláti, quem Pius Papa Duodécimus 
 Sanctórum fastis adscrípsit.
\switchcolumn
\selectlanguage{english}
In the village of Ranft, near Sachseln in Switzerland, St. Nicholas of Flue, 
 a family man who became an anchoret, famed for his most ardent penitence and 
 contempt for the world, and known by the Swiss as the father of the 
 fatherland. He was numbered among the saints by Pope Pius XII.
\switchcolumn*
\selectlanguage{latin}
\end{paracol}


% ---- martyrology/mart03/mart0322.htm
\needspace{10\baselineskip}
\begin{paracol}{2}
\selectlanguage{latin}
\begin{center}{\color{gregoriocolor} Undécimo Kaléndas Aprílis. 
 Luna\dots\ }\end{center}
\switchcolumn
\selectlanguage{english}
\begin{center}{\color{gregoriocolor} The 22\textsuperscript{nd} Day of March. 
 The\dots\ Day of the Moon.}\end{center}
\end{paracol}

\noindent\begin{tabularx}{\linewidth}{*{19}{>{\centering\arraybackslash}X}}
 \textcolor{gregoriocolor}{a} & \textcolor{gregoriocolor}{b} & \textcolor{gregoriocolor}{c} & \textcolor{gregoriocolor}{d} & \textcolor{gregoriocolor}{e} & \textcolor{gregoriocolor}{f} & \textcolor{gregoriocolor}{g} & \textcolor{gregoriocolor}{h} & \textcolor{gregoriocolor}{i} & \textcolor{gregoriocolor}{k} & \textcolor{gregoriocolor}{l} & \textcolor{gregoriocolor}{m} & \textcolor{gregoriocolor}{n} & \textcolor{gregoriocolor}{p} & \textcolor{gregoriocolor}{q} & \textcolor{gregoriocolor}{r} & \textcolor{gregoriocolor}{s} & \textcolor{gregoriocolor}{t} & \textcolor{gregoriocolor}{u} \\
 23 & 24 & 25 & 26 & 27 & 28 & 29 & 30 & 1 & 2 & 3 & 4 & 5 & 6 & 7 & 8 & 9 & 10 & 11 \\
\end{tabularx}
\vspace{0.5\baselineskip}
\noindent\begin{tabularx}{\linewidth}{*{12}{>{\centering\arraybackslash}X}}
 \textcolor{gregoriocolor}{A} & \textcolor{gregoriocolor}{B} & \textcolor{gregoriocolor}{C} & \textcolor{gregoriocolor}{D} & \textcolor{gregoriocolor}{E} & F & \textcolor{gregoriocolor}{F} & \textcolor{gregoriocolor}{G} & \textcolor{gregoriocolor}{H} & \textcolor{gregoriocolor}{M} & \textcolor{gregoriocolor}{N} & \textcolor{gregoriocolor}{P} \\
 12 & 13 & 14 & 15 & 16 & 17 & 17 & 18 & 19 & 20 & 21 & 22 \\
\end{tabularx}

\begin{paracol}{2}
\selectlanguage{latin}
\lettrine[lines=2]{N}{arbóne,} in Gállia, natális sancti Pauli Epíscopi, Apostolórum discípuli, 
 quem tradunt fuísse Sérgium Paulum Procónsulem. Hic, a beáto Apóstolo 
 Paulo baptizátus, et ab eo, cum in Hispániam pérgeret, apud Narbónem 
 relíctus, ibídem Episcopáli dignitáte donátus est; ibíque, prædicatiónis 
 offício non ségniter expléto, clarus miráculis migrávit in cælum.
\switchcolumn
\selectlanguage{english}
\lettrine[lines=2]{A}{t} Narbonne in France, the birthday of the bishop St. Paul, a disciple of 
 the apostles. He is said to have been the proconsul Sergius Paulus, 
 who was baptized by the blessed apostle Paul, and left at Narbonne, where he 
 was raised to the episcopal dignity when the apostle went to Spain. 
 Having zealously discharged the office of preaching and having performed 
 miracles, he departed to heaven.
\switchcolumn*
\selectlanguage{latin}
Tarracínæ, in Campánia, sancti Epaphrodíti, 
 Apostolórum discípuli, qui a beáto Petro Apóstolo Epíscopus illíus civitátis 
 ordinátus fuit.
\switchcolumn
\selectlanguage{english}
At Terracina, St. Epaphroditus, a disciple of the apostles, who was 
 consecrated bishop of that city by the blessed apostle Peter.
\switchcolumn*
\selectlanguage{latin}
Ancyræ, in Galátia, sancti Basilíi, Presbyteri 
 et Mártyris, qui sub Juliáno Apóstata, gravíssimis cruciátibus afféctus, 
 ánimam Deo réddidit.
\switchcolumn
\selectlanguage{english}
At Ancyra, under Julian the Apostate, St. Basil, priest and martyr, who gave 
 up his soul to God after having endured grievous torments.
\switchcolumn*
\selectlanguage{latin}
Carthágine sancti Octaviáni Archidiáconi, et multórum míllium Mártyrum, qui, 
 ob fidem cathólicam, a Wándalis cæsi sunt.
\switchcolumn
\selectlanguage{english}
At Carthage, the archdeacon St. Octavian, and many thousands of martyrs, who 
 were slain by the Vandals for the Catholic faith.
\switchcolumn*
\selectlanguage{latin}
In Africa sanctórum Mártyrum Saturníni et aliórum novem.
\switchcolumn
\selectlanguage{english}
In Africa, the holy martyrs Saturninus and nine others.
\switchcolumn*
\selectlanguage{latin}
In Galátia natális sanctárum Mártyrum Callinícæ 
 et Basilíssæ.
\switchcolumn
\selectlanguage{english}
In Galatia, the birthday of the holy martyrs Callinica and Basilissa.
\switchcolumn*
\selectlanguage{latin}
Romæ sancti Zacharíæ Papæ, qui Dei Ecclésiam 
 summa vigilántia gubernávit, et clarus méritis quiévit in pace.
\switchcolumn
\selectlanguage{english}
At Rome, the birthday of Pope St. Zachary, who governed the Church of God 
 with vigilance, and at last, renowned for miracles, rested in peace.
\switchcolumn*
\selectlanguage{latin}
Carthágine sancti Deográtias, Epíscopi Carthaginénsis, qui plúrimos, a 
 Wándalis captívos ex Urbe ductos, redémit, aliísque sanctis opéribus 
 célebris quiévit in Dómino.
\switchcolumn
\selectlanguage{english}
At Carthage, St. Deogratias, bishop of Carthage, who ransomed many captives 
 taken from that city by the Vandals, and who performed many other good 
 works, after which he went to rest in the Lord.
\switchcolumn*
\selectlanguage{latin}
Auximi, in Picéno, sancti Benvenúti Epíscopi.
\switchcolumn
\selectlanguage{english}
At Osimo, in Piceno, the bishop St. Benvenuto.
\switchcolumn*
\selectlanguage{latin}
Romæ sanctæ Leæ Víduæ, cujus virtútes et 
 tránsitum ad Deum sanctus Hierónymus scribit.
\switchcolumn
\selectlanguage{english}
At Rome, the widow St. Lea, whose virtues and happy death are related by St. 
 Jerome.
\switchcolumn*
\selectlanguage{latin}
\end{paracol}


% ---- martyrology/mart03/mart0323.htm
\needspace{10\baselineskip}
\begin{paracol}{2}
\selectlanguage{latin}
\begin{center}{\color{gregoriocolor} Décimo Kaléndas Aprílis. 
 Luna\dots\ }\end{center}
\switchcolumn
\selectlanguage{english}
\begin{center}{\color{gregoriocolor} The 23\textsuperscript{rd} Day of March. 
 The\dots\ Day of the Moon.}\end{center}
\end{paracol}

\noindent\begin{tabularx}{\linewidth}{*{19}{>{\centering\arraybackslash}X}}
 \textcolor{gregoriocolor}{a} & \textcolor{gregoriocolor}{b} & \textcolor{gregoriocolor}{c} & \textcolor{gregoriocolor}{d} & \textcolor{gregoriocolor}{e} & \textcolor{gregoriocolor}{f} & \textcolor{gregoriocolor}{g} & \textcolor{gregoriocolor}{h} & \textcolor{gregoriocolor}{i} & \textcolor{gregoriocolor}{k} & \textcolor{gregoriocolor}{l} & \textcolor{gregoriocolor}{m} & \textcolor{gregoriocolor}{n} & \textcolor{gregoriocolor}{p} & \textcolor{gregoriocolor}{q} & \textcolor{gregoriocolor}{r} & \textcolor{gregoriocolor}{s} & \textcolor{gregoriocolor}{t} & \textcolor{gregoriocolor}{u} \\
 24 & 25 & 26 & 27 & 28 & 29 & 30 & 1 & 2 & 3 & 4 & 5 & 6 & 7 & 8 & 9 & 10 & 11 & 12 \\
\end{tabularx}
\vspace{0.5\baselineskip}
\noindent\begin{tabularx}{\linewidth}{*{12}{>{\centering\arraybackslash}X}}
 \textcolor{gregoriocolor}{A} & \textcolor{gregoriocolor}{B} & \textcolor{gregoriocolor}{C} & \textcolor{gregoriocolor}{D} & \textcolor{gregoriocolor}{E} & F & \textcolor{gregoriocolor}{F} & \textcolor{gregoriocolor}{G} & \textcolor{gregoriocolor}{H} & \textcolor{gregoriocolor}{M} & \textcolor{gregoriocolor}{N} & \textcolor{gregoriocolor}{P} \\
 13 & 14 & 15 & 16 & 17 & 18 & 18 & 19 & 20 & 21 & 22 & 23 \\
\end{tabularx}

\begin{paracol}{2}
\selectlanguage{latin}
\lettrine[lines=2]{I}{n} Africa sanctórum Mártyrum Victoriáni, Procónsulis Cartháginis, et duórum 
 germanórum, Aquisregénsium; item Fruméntii et altérius Fruméntii, mercatórum. 
 Hi omnes, in persecutióne Wandálica (ut scribit Victor, Africánus Epíscopus), 
 sub Ariáno Rege Hunneríco, pro constántia 
 cathólicæ confessiónis, immaníssimis supplíciis cruciáti, egrégie coronáti 
 sunt.
\switchcolumn
\selectlanguage{english}
\lettrine[lines=2]{I}{n} Africa, the holy martyrs Victorian, proconsul of Carthage, and two 
 brothers from Aquaregia. Also two merchants, both named Frementius, 
 who (as Bishop Victor Africanus relates) were subjected to the most 
 atrocious torments for their courageous confession of the Catholic faith, 
 and who were gloriously crowned martyrs under the Arian king Hunneric, 
 during the persecution of the Vandals.
\switchcolumn*
\selectlanguage{latin}
Item in Africa sancti Fidélis Mártyris.
\switchcolumn
\selectlanguage{english}
Also in Africa, St. Fidelis, martyr.
\switchcolumn*
\selectlanguage{latin}
Ibídem sancti Felícis et aliórum vigínti Mártyrum.
\switchcolumn
\selectlanguage{english}
In the same place, St. Felix and twenty other martyrs.
\switchcolumn*
\selectlanguage{latin}
Cæsaréæ, in Palæstína, sanctórum Mártyrum 
 Nicónis et aliórum nonagínta novem.
\switchcolumn
\selectlanguage{english}
At Caesarea in Palestine, the holy martyrs Nicon and ninety-nine others.
\switchcolumn*
\selectlanguage{latin}
Item corónæ sanctórum Mártyrum Domítii, Pelágiæ, 
 Aquilæ, Epárchii et Theodósiæ.
\switchcolumn
\selectlanguage{english}
Likewise, the crowning of the holy martyrs Domitius, Pelagia, Aquila, 
 Eparchius, and Theodosia.
\switchcolumn*
\selectlanguage{latin}
Limæ, in Perúvia, sancti Turíbii Epíscopi, 
 cujus virtúte fides et disciplína ecclesiástica per Amerícam diffúsæ sunt.
\switchcolumn
\selectlanguage{english}
At Lima in Peru, Archbishop St. Turibius, through whose labours both faith 
 and ecclesiastical discipline were spread through America.
\switchcolumn*
\selectlanguage{latin}
Antiochíæ sancti Theodúli Presbyteri.
\switchcolumn
\selectlanguage{english}
At Antioch, the priest St. Theodulus.
\switchcolumn*
\selectlanguage{latin}
Barcinóne, in Hispánia, sancti Joséphi Oriol Presbyteri, Ecclésiæ 
 sanctæ Maríæ Regum Beneficiárii, omnígena virtúte ac præsértim córporis 
 afflictatióne, paupertátis cultu atque in egénos et infírmos caritáte 
 célebris; quem, in vita et post mortem miráculis gloriósum, Pius Papa 
 Décimus in Sanctórum númerum recénsuit.
\switchcolumn
\selectlanguage{english}
At Barcelona in Spain, the priest St. Joseph Oriol, pastor of the church of 
 St. Mary of the Kings, famous for every virtue, especially mortification of 
 the body, his rule of poverty, and his love towards the poor and the sick. 
 Because he was known for his miracles both in life and after death, Pope 
 Pius X placed his name in the number of the saints.
\switchcolumn*
\selectlanguage{latin}
Cæsaréæ sancti Juliáni Confessóris.
\switchcolumn
\selectlanguage{english}
At Caesarea, St. Julian, confessor.
\switchcolumn*
\selectlanguage{latin}
In Campánia sancti Benedícti Mónachi, qui, a Gothis in ardénti clíbano 
 inclúsus, sequénti die invéntus est illæsus.
\switchcolumn
\selectlanguage{english}
In Campania, St. Benedict, monk, who was shut up in a burning furnace by the 
 Goths, but who was found uninjured the next day.
\switchcolumn*
\selectlanguage{latin}
\end{paracol}


% ---- martyrology/mart03/mart0324.htm
\needspace{10\baselineskip}
\begin{paracol}{2}
\selectlanguage{latin}
\begin{center}{\color{gregoriocolor} Nono Kaléndas Aprílis. 
 Luna\dots\ }\end{center}
\switchcolumn
\selectlanguage{english}
\begin{center}{\color{gregoriocolor} The 24\textsuperscript{th} Day of March. 
 The\dots\ Day of the Moon.}\end{center}
\end{paracol}

\noindent\begin{tabularx}{\linewidth}{*{19}{>{\centering\arraybackslash}X}}
 \textcolor{gregoriocolor}{a} & \textcolor{gregoriocolor}{b} & \textcolor{gregoriocolor}{c} & \textcolor{gregoriocolor}{d} & \textcolor{gregoriocolor}{e} & \textcolor{gregoriocolor}{f} & \textcolor{gregoriocolor}{g} & \textcolor{gregoriocolor}{h} & \textcolor{gregoriocolor}{i} & \textcolor{gregoriocolor}{k} & \textcolor{gregoriocolor}{l} & \textcolor{gregoriocolor}{m} & \textcolor{gregoriocolor}{n} & \textcolor{gregoriocolor}{p} & \textcolor{gregoriocolor}{q} & \textcolor{gregoriocolor}{r} & \textcolor{gregoriocolor}{s} & \textcolor{gregoriocolor}{t} & \textcolor{gregoriocolor}{u} \\
 25 & 26 & 27 & 28 & 29 & 30 & 1 & 2 & 3 & 4 & 5 & 6 & 7 & 8 & 9 & 10 & 11 & 12 & 13 \\
\end{tabularx}
\vspace{0.5\baselineskip}
\noindent\begin{tabularx}{\linewidth}{*{12}{>{\centering\arraybackslash}X}}
 \textcolor{gregoriocolor}{A} & \textcolor{gregoriocolor}{B} & \textcolor{gregoriocolor}{C} & \textcolor{gregoriocolor}{D} & \textcolor{gregoriocolor}{E} & F & \textcolor{gregoriocolor}{F} & \textcolor{gregoriocolor}{G} & \textcolor{gregoriocolor}{H} & \textcolor{gregoriocolor}{M} & \textcolor{gregoriocolor}{N} & \textcolor{gregoriocolor}{P} \\
 14 & 15 & 16 & 17 & 18 & 19 & 19 & 20 & 21 & 22 & 23 & 24 \\
\end{tabularx}

\begin{paracol}{2}
\selectlanguage{latin}
\lettrine[lines=2]{F}{estum} sancti Gabriélis Archángeli, qui ad annuntiándum Incarnatiónis divíni 
 Verbi mystérium a Deo missus est.
\switchcolumn
\selectlanguage{english}
\lettrine[lines=2]{T}{he} Feast of St. Gabriel Archangel, who was sent by God to announce the 
 Incarnation of the Divine Word.
\switchcolumn*
\selectlanguage{latin}
Romæ sancti Epigménii Presbyteri, qui, in 
 persecutióne Diocletiáni, sub Túrpio Júdice, gládio cæsus, martyrium 
 consummávit.
\switchcolumn
\selectlanguage{english}
At Rome, the priest St. Epigmenius, who completed his martyrdom by the sword 
 in the persecution of Diocletian, under the judge Turpius.
\switchcolumn*
\selectlanguage{latin}
Ibídem pássio beáti Pigménii Presbyteri, qui, sub Juliáno Apóstata, pro fide 
 Christi, præcipitátus in Tíberim, necátus est.
\switchcolumn
\selectlanguage{english}
Also at Rome, in the time of Julian the Apostate, the passion of blessed 
 Pigmenius, a priest, who was killed for the faith of Christ by being drowned 
 in the Tiber.
\switchcolumn*
\selectlanguage{latin}
Item Romæ sanctórum Mártyrum Marci et Timóthei, 
 qui martyrio coronáti sunt sub Antoníno Imperatóre.
\switchcolumn
\selectlanguage{english}
At Rome, the holy martyrs Mark and Timothy, who were crowned with martyrdom 
 under Emperor Antoninus.
\switchcolumn*
\selectlanguage{latin}
Cæsaréæ, in Palæstína, natális sanctórum 
 Mártyrum Timolái, Dionysii, Páusidis, Rómuli, Alexándri, altérius Alexándri, 
 Agápii et altérius Dionysii; qui, in persecutióne Diocletiáni, sub Urbáno 
 Præside, secúris ictu percússi, vitæ corónas meruérunt.
\switchcolumn
\selectlanguage{english}
At Caesarea in Palestine, the birthday of the holy martyrs Timolaus, Denis, 
 Pausides, Romulus, Alexander, another Alexander, Agapius, and another Denis. 
 They merited the crown of life by being beheaded in the persecution of 
 Diocletian under the governor Urban.
\switchcolumn*
\selectlanguage{latin}
In Mauritánia item natális sanctórum fratrum Rómuli et Secúndi, qui pro 
 Christi fide passi sunt.
\switchcolumn
\selectlanguage{english}
In Morocco, the birthday of the saintly brothers Romulus and Secundus, who 
 suffered for the faith of Christ.
\switchcolumn*
\selectlanguage{latin}
Tridénti pássio sancti Simeónis púeri, a Judǽis 
 sævíssime trucidáti, qui multis póstea miráculis coruscávit.
\switchcolumn
\selectlanguage{english}
At Trent, the martyrdom of the boy St. Simeon, who was barbarously murdered 
 by the Jews, but who was afterwards glorified by many miracles.
\switchcolumn*
\selectlanguage{latin}
Synnadæ, in Phrygia, sancti Agapíti Epíscopi.
\switchcolumn
\selectlanguage{english}
At Synnadas in Phrygia, Bishop St. Agapitus.
\switchcolumn*
\selectlanguage{latin}
Bríxiæ sancti Latíni Epíscopi.
\switchcolumn
\selectlanguage{english}
At Brescia, the bishop St. Latinus.
\switchcolumn*
\selectlanguage{latin}
In Syria sancti Seléuci Confessóris.
\switchcolumn
\selectlanguage{english}
In Syria, St. Seleucus, confessor.
\switchcolumn*
\selectlanguage{latin}
In Suécia sanctæ Catharínæ Vírginis, quæ fuit 
 fília sanctæ Birgíttæ.
\switchcolumn
\selectlanguage{english}
In Sweden, the virgin St. Catherine, daughter of St. Bridget.
\switchcolumn*
\selectlanguage{latin}
\end{paracol}


% ---- martyrology/mart03/mart0325.htm
\needspace{10\baselineskip}
\begin{paracol}{2}
\selectlanguage{latin}
\begin{center}{\color{gregoriocolor} Octávo Kaléndas Aprílis. 
 Luna\dots\ }\end{center}
\switchcolumn
\selectlanguage{english}
\begin{center}{\color{gregoriocolor} The 25\textsuperscript{th} Day of March. 
 The\dots\ Day of the Moon.}\end{center}
\end{paracol}

\noindent\begin{tabularx}{\linewidth}{*{19}{>{\centering\arraybackslash}X}}
 \textcolor{gregoriocolor}{a} & \textcolor{gregoriocolor}{b} & \textcolor{gregoriocolor}{c} & \textcolor{gregoriocolor}{d} & \textcolor{gregoriocolor}{e} & \textcolor{gregoriocolor}{f} & \textcolor{gregoriocolor}{g} & \textcolor{gregoriocolor}{h} & \textcolor{gregoriocolor}{i} & \textcolor{gregoriocolor}{k} & \textcolor{gregoriocolor}{l} & \textcolor{gregoriocolor}{m} & \textcolor{gregoriocolor}{n} & \textcolor{gregoriocolor}{p} & \textcolor{gregoriocolor}{q} & \textcolor{gregoriocolor}{r} & \textcolor{gregoriocolor}{s} & \textcolor{gregoriocolor}{t} & \textcolor{gregoriocolor}{u} \\
 26 & 27 & 28 & 29 & 30 & 1 & 2 & 3 & 4 & 5 & 6 & 7 & 8 & 9 & 10 & 11 & 12 & 13 & 14 \\
\end{tabularx}
\vspace{0.5\baselineskip}
\noindent\begin{tabularx}{\linewidth}{*{12}{>{\centering\arraybackslash}X}}
 \textcolor{gregoriocolor}{A} & \textcolor{gregoriocolor}{B} & \textcolor{gregoriocolor}{C} & \textcolor{gregoriocolor}{D} & \textcolor{gregoriocolor}{E} & F & \textcolor{gregoriocolor}{F} & \textcolor{gregoriocolor}{G} & \textcolor{gregoriocolor}{H} & \textcolor{gregoriocolor}{M} & \textcolor{gregoriocolor}{N} & \textcolor{gregoriocolor}{P} \\
 15 & 16 & 17 & 18 & 19 & 20 & 20 & 21 & 22 & 23 & 24 & 25 \\
\end{tabularx}

\begin{paracol}{2}
\selectlanguage{latin}
\lettrine[lines=2]{A}{nnuntiátio} beatíssimæ Vírginis Genitrícis Dei 
 Maríæ.
\switchcolumn
\selectlanguage{english}
\lettrine[lines=2]{T}{he} Annunciation of the Most Blessed Virgin Mary, Mother of God.
\switchcolumn*
\selectlanguage{latin}
Hierosólymis commemorátio sancti Latrónis, qui, in cruce Christum conféssus, 
 ab eo méruit audíre: « Hódie mecum eris in 
 paradíso ».
\switchcolumn
\selectlanguage{english}
At Jerusalem, the commemoration of the good thief who confessed Christ on 
 the cross, and who deserved to hear from him these words: ``This day shalt 
 thou be with me in paradise.''
\switchcolumn*
\selectlanguage{latin}
Romæ sancti Quiríni Mártyris, qui, sub Cláudio 
 Imperatóre, post facultátum amissiónem, post cárceris squalórem, post 
 multórum vérberum afflictiónem, gládio interféctus est et in Tíberim 
 projéctus; quem Christiáni, cum in ínsula Lycaónia (quæ póstea sancti 
 Bartholomæi dicta est) inveníssent, in cœmetério Pontiáni condidérunt.
\switchcolumn
\selectlanguage{english}
At Rome, St. Quirinus, martyr, who after losing his possessions, suffering 
 imprisonment in a dark dungeon, and being cruelly scourged, was put to death 
 with the sword, and thrown into the Tiber. The Christians found his 
 body on the island of Lycaonia (which was thereafter called St. 
 Bartholomew's), and buried it in the Pontian cemetery.
\switchcolumn*
\selectlanguage{latin}
Item Romæ sanctórum ducentórum sexagínta duórum 
 Mártyrum.
\switchcolumn
\selectlanguage{english}
Also at Rome, two hundred and sixty-two holy martyrs.
\switchcolumn*
\selectlanguage{latin}
Sírmii pássio sancti Irenǽi, Epíscopi et 
 Mártyris; qui, témpore Maximiáni Imperatóris, sub Præside Probo, 
 primum torméntis acérrimis vexátus, deínde diébus plúrimis cruciátus in cárcere, novíssime, abscísso cápite, consummátus est.
\switchcolumn
\selectlanguage{english}
At Sirmium, the martyrdom of St. Irenaeus, bishop. In the time of 
 Emperor Maximian, under the governor Probus, after undergoing bitter 
 torments and a painful imprisonment for may days, he was beheaded.
\switchcolumn*
\selectlanguage{latin}
Nicomedíæ sanctæ Dulæ, cujúsdam mílitis ancíllæ, 
 quæ, ob castitátem servándam occísa, martyrii corónam proméruit.
\switchcolumn
\selectlanguage{english}
At Nicomedia, St. Dula, the servant of a soldier, who was killed for the 
 preservation of her chastity, and deserved the crown of martyrdom.
\switchcolumn*
\selectlanguage{latin}
Laodicéæ, ad Líbanum, sancti Pelágii Epíscopi, 
 qui, ob fidem cathólicam, témpore Valéntis, exsílium et ália passus est; ac 
 tandem, in sedem suam restitútus, quiévit in Dómino.
\switchcolumn
\selectlanguage{english}
At Laodicea, St. Pelagius, bishop, who after having endured exile and other 
 afflictions for the Catholic faith under Valens, rested in the Lord.
\switchcolumn*
\selectlanguage{latin}
In Antro, ínsula Lígeris flúminis, sancti 
 Hermelándi Abbátis, cujus gloriósa conversátio insígni miraculórum præcónio 
 commendátur.
\switchcolumn
\selectlanguage{english}
At Indre, an island in the Loire, Abbot St. Hermeland, whose glorious life 
 was commended by outstanding miracles.
\switchcolumn*
\selectlanguage{latin}
Pistórii, in Túscia, sanctórum Confessórum Baróntii et Desidérii.
\switchcolumn
\selectlanguage{english}
At Pistoia, the holy confessors Barontius and Desiderius.
\switchcolumn*
\selectlanguage{latin}
Faliscodúni sanctæ 
 Lúciæ Filippíni, Fundatrícis Institúti Magistrárum Piárum ab ejus cognómine 
 nuncupatárum, de Christiána puellárum et mulíerum, præsértim páuperum, 
 eruditióne óptime méritæ, quam Pius Papa Undécimus inter sanctas Vírgines 
 rétulit.
\switchcolumn
\selectlanguage{english}
At Montefiascone, St. Lucia Filippini, founder of the Institute of Pious 
 Teachers, from whose surname they are known as Filippines. 
 Having merited greatly by the Christian education of girls and women, 
 especially of the poor, Pope Pius XI enrolled her among the holy virgins.
\switchcolumn*
\selectlanguage{latin}
\end{paracol}


% ---- martyrology/mart03/mart0326.htm
\needspace{10\baselineskip}
\begin{paracol}{2}
\selectlanguage{latin}
\begin{center}{\color{gregoriocolor} Séptimo Kaléndas Aprílis. 
 Luna\dots\ }\end{center}
\switchcolumn
\selectlanguage{english}
\begin{center}{\color{gregoriocolor} The 26\textsuperscript{th} Day of March. 
 The\dots\ Day of the Moon.}\end{center}
\end{paracol}

\noindent\begin{tabularx}{\linewidth}{*{19}{>{\centering\arraybackslash}X}}
 \textcolor{gregoriocolor}{a} & \textcolor{gregoriocolor}{b} & \textcolor{gregoriocolor}{c} & \textcolor{gregoriocolor}{d} & \textcolor{gregoriocolor}{e} & \textcolor{gregoriocolor}{f} & \textcolor{gregoriocolor}{g} & \textcolor{gregoriocolor}{h} & \textcolor{gregoriocolor}{i} & \textcolor{gregoriocolor}{k} & \textcolor{gregoriocolor}{l} & \textcolor{gregoriocolor}{m} & \textcolor{gregoriocolor}{n} & \textcolor{gregoriocolor}{p} & \textcolor{gregoriocolor}{q} & \textcolor{gregoriocolor}{r} & \textcolor{gregoriocolor}{s} & \textcolor{gregoriocolor}{t} & \textcolor{gregoriocolor}{u} \\
 27 & 28 & 29 & 30 & 1 & 2 & 3 & 4 & 5 & 6 & 7 & 8 & 9 & 10 & 11 & 12 & 13 & 14 & 15 \\
\end{tabularx}
\vspace{0.5\baselineskip}
\noindent\begin{tabularx}{\linewidth}{*{12}{>{\centering\arraybackslash}X}}
 \textcolor{gregoriocolor}{A} & \textcolor{gregoriocolor}{B} & \textcolor{gregoriocolor}{C} & \textcolor{gregoriocolor}{D} & \textcolor{gregoriocolor}{E} & F & \textcolor{gregoriocolor}{F} & \textcolor{gregoriocolor}{G} & \textcolor{gregoriocolor}{H} & \textcolor{gregoriocolor}{M} & \textcolor{gregoriocolor}{N} & \textcolor{gregoriocolor}{P} \\
 16 & 17 & 18 & 19 & 20 & 21 & 21 & 22 & 23 & 24 & 25 & 26 \\
\end{tabularx}

\begin{paracol}{2}
\selectlanguage{latin}
\lettrine[lines=2]{R}{omæ,} via Lavicána, sancti Cástuli Mártyris, 
 qui, cum esset zetárius Palátii et hospes Sanctórum, a persecutóribus tértio 
 appénsus, tértio audítus, et, in confessióne Dómini persevérans, missus est 
 in fóveam, ac, dimíssa super eum massa arenária, martyrio coronátus est.
\switchcolumn
\selectlanguage{english}
\lettrine[lines=2]{A}{t} Rome, on the Via Lavicana, St. Castulus, martyr, chamberlain in the 
 palace of the emperor. For harbouring Christians, he was three times 
 suspended by the hands, three times cited before the tribunals. As he 
 persevered in the confession of the Lord, he was thrown into a pit, covered 
 with a mass of sand, and thus obtained the crown of martyrdom.
\switchcolumn*
\selectlanguage{latin}
Item Romæ corónæ sanctórum Mártyrum Petri, 
 Marciáni, Jovíni, Theclæ, Cassiáni et aliórum.
\switchcolumn
\selectlanguage{english}
Also at Rome, the crowning of the holy martyrs Peter, Marcian, Jovinus, 
 Thecla, Cassian, and others.
\switchcolumn*
\selectlanguage{latin}
Pentápoli, in Libya, natális sanctórum 
 Mártyrum Theodóri Epíscopi, Irenǽi Diáconi, 
 Serapiónis et Ammónii Lectórum.
\switchcolumn
\selectlanguage{english}
At Pentopolis in Libya, the birthday of the holy martyrs Theodore, bishop, 
 the deacon Irenæus, and the lectors Serapion and Ammonius.
\switchcolumn*
\selectlanguage{latin}
Sírmii sanctórum Mártyrum Montáni Presbyteri, et Máximæ, 
 qui, ob Christi fidem, in flumen demérsi sunt.
\switchcolumn
\selectlanguage{english}
At Sirmium, the holy martyrs Montanus, priest, and Maxima, who were drowned 
 in a river for the faith of Christ.
\switchcolumn*
\selectlanguage{latin}
Item sanctórum Mártyrum Quadráti, Theodósii, Emmanuélis et aliórum 
 quadragínta.
\switchcolumn
\selectlanguage{english}
Likewise, the holy martyrs Quadratus, Theodosius, Emmanuel, and forty 
 others.
\switchcolumn*
\selectlanguage{latin}
Alexandríæ sanctórum Mártyrum Eutychii et 
 aliórum; qui, Constántii témpore, sub Ariáno Epíscopo Geórgio, pro fide 
 cathólica gládio cæsi sunt.
\switchcolumn
\selectlanguage{english}
At Alexandria, the holy martyrs Eutychius and others, who died by the sword 
 for the Catholic faith, in the time of Constantine, under the Arian bishop 
 George.
\switchcolumn*
\selectlanguage{latin}
Eódem die sancti Ludgéri, Epíscopi Monasteriénsis, qui Saxónibus Evangélium 
 prædicávit.
\switchcolumn
\selectlanguage{english}
The same day, St. Ludger, bishop of Munster, who preached the Gospel to the 
 Saxons.
\switchcolumn*
\selectlanguage{latin}
Cæsaraugústæ, in Hispánia, sancti Bráulii, 
 Epíscopi et Confessóris.
\switchcolumn
\selectlanguage{english}
At Saragossa in Spain, St. Braulio, bishop and confessor.
\switchcolumn*
\selectlanguage{latin}
Tréviris sancti Felícis Epíscopi.
\switchcolumn
\selectlanguage{english}
At Treves, St. Felix, bishop.
\switchcolumn*
\selectlanguage{latin}
\end{paracol}


% ---- martyrology/mart03/mart0327.htm
\needspace{10\baselineskip}
\begin{paracol}{2}
\selectlanguage{latin}
\begin{center}{\color{gregoriocolor} Sexto Kaléndas Aprílis. 
 Luna\dots\ }\end{center}
\switchcolumn
\selectlanguage{english}
\begin{center}{\color{gregoriocolor} The 27\textsuperscript{th} Day of March. 
 The\dots\ Day of the Moon.}\end{center}
\end{paracol}

\noindent\begin{tabularx}{\linewidth}{*{19}{>{\centering\arraybackslash}X}}
 \textcolor{gregoriocolor}{a} & \textcolor{gregoriocolor}{b} & \textcolor{gregoriocolor}{c} & \textcolor{gregoriocolor}{d} & \textcolor{gregoriocolor}{e} & \textcolor{gregoriocolor}{f} & \textcolor{gregoriocolor}{g} & \textcolor{gregoriocolor}{h} & \textcolor{gregoriocolor}{i} & \textcolor{gregoriocolor}{k} & \textcolor{gregoriocolor}{l} & \textcolor{gregoriocolor}{m} & \textcolor{gregoriocolor}{n} & \textcolor{gregoriocolor}{p} & \textcolor{gregoriocolor}{q} & \textcolor{gregoriocolor}{r} & \textcolor{gregoriocolor}{s} & \textcolor{gregoriocolor}{t} & \textcolor{gregoriocolor}{u} \\
 28 & 29 & 30 & 1 & 2 & 3 & 4 & 5 & 6 & 7 & 8 & 9 & 10 & 11 & 12 & 13 & 14 & 15 & 16 \\
\end{tabularx}
\vspace{0.5\baselineskip}
\noindent\begin{tabularx}{\linewidth}{*{12}{>{\centering\arraybackslash}X}}
 \textcolor{gregoriocolor}{A} & \textcolor{gregoriocolor}{B} & \textcolor{gregoriocolor}{C} & \textcolor{gregoriocolor}{D} & \textcolor{gregoriocolor}{E} & F & \textcolor{gregoriocolor}{F} & \textcolor{gregoriocolor}{G} & \textcolor{gregoriocolor}{H} & \textcolor{gregoriocolor}{M} & \textcolor{gregoriocolor}{N} & \textcolor{gregoriocolor}{P} \\
 17 & 18 & 19 & 20 & 21 & 22 & 22 & 23 & 24 & 25 & 26 & 27 \\
\end{tabularx}

\begin{paracol}{2}
\selectlanguage{latin}
\lettrine[lines=2]{S}{ancti} Joánnis Damascéni, Presbyteri, Confessóris et Ecclésiæ 
 Doctoris, cujus dies natális ágitur prídie Nonas Maji.
\switchcolumn
\selectlanguage{english}
\lettrine[lines=2]{S}{t.} John Damascene, priest, confessor, and doctor of the Church, whose 
 birthday is commemorated on the 6th of May.
\switchcolumn*
\selectlanguage{latin}
Drizíparæ, 
 in Pannónia, sancti Alexándri mílitis, qui, sub Maximiáno Imperatóre, post 
 multos pro Christo agónes superátos múltaque mirácula édita, cápitis 
 abscissióne martyrium complévit.
\switchcolumn
\selectlanguage{english}
At Drizipara in Hungary, St. Alexander, soldier, in the time of Emperor 
 Maximian. Having overcome many torments for the sake of Christ, and 
 performing many miracles, his martyrdom was completed by beheading.
\switchcolumn*
\selectlanguage{latin}
In Illyrico sanctórum Philéti Senatóris, Lydiæ 
 uxóris, et filiórum Macédonis et Theoprépii, itémque Amphilóchii Ducis, et 
 Crónidæ Commentariénsis; qui, pro Christi confessióne, torméntis plúribus 
 superátis, corónam glóriæ sunt adépti.
\switchcolumn
\selectlanguage{english}
In Illyria, the Saints Philetus, senator, his wife Lydia, and their sons 
 Macedon and Theoprepides; also Amphilochius, an officer in the army, and 
 Chronides, a notary, who were put to death for the confession of Christ 
 after suffering many things.
\switchcolumn*
\selectlanguage{latin}
In Pérside sanctórum Mártyrum Zanítæ, Lázari, 
 Marótæ, Narsétis et aliórum quinque, qui sub Rege Persárum Sápore, sævíssime 
 trucidáti, martyrii palmam meruérunt.
\switchcolumn
\selectlanguage{english}
In Persia, in the reign of King Sapor, the holy martyrs Zanitas, Lazarus, 
 Marotas, Narses, and five others, who were barbarously slain, having merited 
 the martyr's palm.
\switchcolumn*
\selectlanguage{latin}
Salisbúrgi, in Nórico, sancti Rupérti, Epíscopi 
 et Confessóris, qui apud Bávaros et Nóricos Evangélium mirífice propagávit.
\switchcolumn
\selectlanguage{english}
At Salzburg in Austria, St. Rupert, bishop and confessor, who spread the 
 Gospel extensively in Bavaria and Austria.
\switchcolumn*
\selectlanguage{latin}
In Ægypto sancti Joánnis Eremítæ, magnæ 
 sanctitátis viri, qui, inter cétera virtútum insígnia, étiam prophético 
 spíritu plenus, Theodósio Imperatóri victórias de tyránnis Máximo et Eugénio 
 prædíxit.
\switchcolumn
\selectlanguage{english}
In Egypt, the hermit St. John, a man of great sanctity, who, among other 
 virtues, was filled with the spirit of prophecy, and predicted to Emperor 
 Theodosius his victories over the tyrants Maximus and Eugene.
\switchcolumn*
\selectlanguage{latin}
\end{paracol}


% ---- martyrology/mart03/mart0328.htm
\needspace{10\baselineskip}
\begin{paracol}{2}
\selectlanguage{latin}
\begin{center}{\color{gregoriocolor} Quinto Kaléndas Aprílis. 
 Luna\dots\ }\end{center}
\switchcolumn
\selectlanguage{english}
\begin{center}{\color{gregoriocolor} The 28\textsuperscript{th} Day of March. 
 The\dots\ Day of the Moon.}\end{center}
\end{paracol}

\noindent\begin{tabularx}{\linewidth}{*{19}{>{\centering\arraybackslash}X}}
 \textcolor{gregoriocolor}{a} & \textcolor{gregoriocolor}{b} & \textcolor{gregoriocolor}{c} & \textcolor{gregoriocolor}{d} & \textcolor{gregoriocolor}{e} & \textcolor{gregoriocolor}{f} & \textcolor{gregoriocolor}{g} & \textcolor{gregoriocolor}{h} & \textcolor{gregoriocolor}{i} & \textcolor{gregoriocolor}{k} & \textcolor{gregoriocolor}{l} & \textcolor{gregoriocolor}{m} & \textcolor{gregoriocolor}{n} & \textcolor{gregoriocolor}{p} & \textcolor{gregoriocolor}{q} & \textcolor{gregoriocolor}{r} & \textcolor{gregoriocolor}{s} & \textcolor{gregoriocolor}{t} & \textcolor{gregoriocolor}{u} \\
 29 & 30 & 1 & 2 & 3 & 4 & 5 & 6 & 7 & 8 & 9 & 10 & 11 & 12 & 13 & 14 & 15 & 16 & 17 \\
\end{tabularx}
\vspace{0.5\baselineskip}
\noindent\begin{tabularx}{\linewidth}{*{12}{>{\centering\arraybackslash}X}}
 \textcolor{gregoriocolor}{A} & \textcolor{gregoriocolor}{B} & \textcolor{gregoriocolor}{C} & \textcolor{gregoriocolor}{D} & \textcolor{gregoriocolor}{E} & F & \textcolor{gregoriocolor}{F} & \textcolor{gregoriocolor}{G} & \textcolor{gregoriocolor}{H} & \textcolor{gregoriocolor}{M} & \textcolor{gregoriocolor}{N} & \textcolor{gregoriocolor}{P} \\
 18 & 19 & 20 & 21 & 22 & 23 & 23 & 24 & 25 & 26 & 27 & 28 \\
\end{tabularx}

\begin{paracol}{2}
\selectlanguage{latin}
\lettrine[lines=2]{S}{ancti} Joánnis de Capistráno, Sacerdótis ex Ordine Minórum et Confessóris, 
 cujus memória recólitur décimo Kaléndas Novémbris.
\switchcolumn
\selectlanguage{english}
\lettrine[lines=2]{S}{t.} John Capistrano, confessor, a priest of the Order of Friars Minor, who 
 is mentioned on the 23rd of October.
\switchcolumn*
\selectlanguage{latin}
Cæsaréæ, in Palæstína, natális sanctórum 
 Mártyrum Prisci, Malchi et Alexándri. Hi tres, in persecutióne 
 Valeriáni, cum in suburbáno agéllo supradíctæ urbis habitárent, atque in ea 
 cæléstes martyrii proponeréntur corónæ, ultro Júdicem, divíno fídei calóre 
 succénsi, ádeunt, et cur tantum in sánguinem piórum desævíret, objúrgant; 
 quos ille contínuo, pro Christi nómine, béstiis trádidit devorándos.
\switchcolumn
\selectlanguage{english}
At Caesarea in Palestine, the birthday of the holy martyrs Priscus, Malchus, 
 and Alexander. In the persecution of Valerian, they were living the 
 suburbs of Caesarea, but knowing that in the city the heavenly crown of 
 martyrdom was to be gained, and burning with the divine ardour of faith, 
 they went to the judge of their own accord, rebuked him for shedding in 
 torrents the blood of the faithful, and were immediately condemned to be 
 devoured by beasts for the Name of Christ.
\switchcolumn*
\selectlanguage{latin}
Tarsi, in Cilícia, sanctórum Mártyrum Cástoris 
 et Doróthei.
\switchcolumn
\selectlanguage{english}
At Tarsus in Cilicia, the holy martyrs Castor and Dorotheus.
\switchcolumn*
\selectlanguage{latin}
In Africa sanctórum Mártyrum Rogáti, 
 Succéssi et aliórum séxdecim.
\switchcolumn
\selectlanguage{english}
In Africa, the holy martyrs Rogatus, Successus, and sixteen others.
\switchcolumn*
\selectlanguage{latin}
Apud Núrsiam sancti Spei Abbátis, miræ 
 patiéntiæ viri, cujus ánima (ut refert sanctus Gregórius Papa), cum ex hac 
 vita migráret, in colúmbæ spécie a cunctis frátribus visa est in cælum 
 ascéndere.
\switchcolumn
\selectlanguage{english}
At Norcia, Abbot St. Spes, a man of extraordinary patience, whose soul at 
 its departure from this life (as Pope St. Gregory relates) was seen by all 
 his brethren to ascend to heaven in the shape of a dove.
\switchcolumn*
\selectlanguage{latin}
Cabillóne, in Gálliis, deposítio sancti Gunthrámni, Regis Francórum, qui 
 spiritálibus actiónibus ita se mancipávit, ut, relíctis, sæculi 
 pompis, thesáuros suos lárgiter Ecclésiis et paupéribus erogáret.
\switchcolumn
\selectlanguage{english}
At Chalons in France, the death of St. Guntram, king of the Franks, who 
 devoted himself to exercises of piety, despising the ostentation of the 
 world, and who bestowed his treasures on the Church and the poor.
\switchcolumn*
\selectlanguage{latin}
\end{paracol}


% ---- martyrology/mart03/mart0329.htm
\needspace{10\baselineskip}
\begin{paracol}{2}
\selectlanguage{latin}
\begin{center}{\color{gregoriocolor} Quarto Kaléndas Aprílis. 
 Luna\dots\ }\end{center}
\switchcolumn
\selectlanguage{english}
\begin{center}{\color{gregoriocolor} The 29\textsuperscript{th} Day of March. 
 The\dots\ Day of the Moon.}\end{center}
\end{paracol}

\noindent\begin{tabularx}{\linewidth}{*{19}{>{\centering\arraybackslash}X}}
 \textcolor{gregoriocolor}{a} & \textcolor{gregoriocolor}{b} & \textcolor{gregoriocolor}{c} & \textcolor{gregoriocolor}{d} & \textcolor{gregoriocolor}{e} & \textcolor{gregoriocolor}{f} & \textcolor{gregoriocolor}{g} & \textcolor{gregoriocolor}{h} & \textcolor{gregoriocolor}{i} & \textcolor{gregoriocolor}{k} & \textcolor{gregoriocolor}{l} & \textcolor{gregoriocolor}{m} & \textcolor{gregoriocolor}{n} & \textcolor{gregoriocolor}{p} & \textcolor{gregoriocolor}{q} & \textcolor{gregoriocolor}{r} & \textcolor{gregoriocolor}{s} & \textcolor{gregoriocolor}{t} & \textcolor{gregoriocolor}{u} \\
 30 & 1 & 2 & 3 & 4 & 5 & 6 & 7 & 8 & 9 & 10 & 11 & 12 & 13 & 14 & 15 & 16 & 17 & 18 \\
\end{tabularx}
\vspace{0.5\baselineskip}
\noindent\begin{tabularx}{\linewidth}{*{12}{>{\centering\arraybackslash}X}}
 \textcolor{gregoriocolor}{A} & \textcolor{gregoriocolor}{B} & \textcolor{gregoriocolor}{C} & \textcolor{gregoriocolor}{D} & \textcolor{gregoriocolor}{E} & F & \textcolor{gregoriocolor}{F} & \textcolor{gregoriocolor}{G} & \textcolor{gregoriocolor}{H} & \textcolor{gregoriocolor}{M} & \textcolor{gregoriocolor}{N} & \textcolor{gregoriocolor}{P} \\
 19 & 20 & 21 & 22 & 23 & 24 & 24 & 25 & 26 & 27 & 28 & 29 \\
\end{tabularx}

\begin{paracol}{2}
\selectlanguage{latin}
\lettrine[lines=2]{H}{eliópoli,} apud Líbanum, 
 sancti Cyrílli, Diáconi et Mártyris, cujus jecur, e discísso ventre avúlsum, 
 Gentíles, sub Juliáno Apóstata, feráliter depásti sunt.
\switchcolumn
\selectlanguage{english}
\lettrine[lines=2]{A}{t} Heliopolis in Lebanon, under Julian the Apostate, St. Cyril, deacon and 
 martyr, whose body was opened and his liver taken out by the heathens who 
 devoured it like wild beasts.
\switchcolumn*
\selectlanguage{latin}
In Pérside sanctórum Monachórum et Mártyrum Jonæ 
 et Barachísii fratrum, sub Rege Persárum Sápore. Ex ipsis Jonas, 
 préssus in cóchlea, confráctis óssibus, médius disséctus est; Barachísius 
 autem, opplétis fáucibus pice ardénti suffocátus.
\switchcolumn
\selectlanguage{english}
In Persia, the holy martyrs Jonas and Barachisius, under the Persian king 
 Sapor. Jonas was put under the pressure of a vice, his bones broken, 
 and cut asunder; Barachisius was suffocated by burning pitch being poured 
 into his throat.
\switchcolumn*
\selectlanguage{latin}
Nicomedíæ pássio sanctórum Mártyrum Pastóris, 
 Victoríni et Sociórum.
\switchcolumn
\selectlanguage{english}
At Nicomedia, the passion of the holy martyrs Pastor, Victorinus, and their 
 companions.
\switchcolumn*
\selectlanguage{latin}
In Africa sanctórum Confessórum Armogástis Cómitis, 
 Másculæ archimími, et Satúri, régiæ domus procuratóris; qui, témpore 
 Wandálicæ persecutiónis, sub Rege Ariáno Genséríco, pro confessióne 
 veritátis, multa et grávia perpéssi supplícia atque oppróbria, cursum 
 gloriósi certáminis implevérunt.
\switchcolumn
\selectlanguage{english}
In Africa, under the Arian king Genseric, during the persecution of the 
 Vandals, the holy confessors Armogastes, a count, Mascula, Archimimus, and 
 Saturus, master of the king's household. After enduring many severe 
 torments, as well as insults, for the confession of the truth, they 
 completed their tests with glory.
\switchcolumn*
\selectlanguage{latin}
In urbe Asténsi sancti Secúndi Mártyris.
\switchcolumn
\selectlanguage{english}
In the town of Asti, St. Secundus, martyr.
\switchcolumn*
\selectlanguage{latin}
In monastério Luxoviénsi, in Gállia, deposítio sancti Eustásii Abbátis, qui 
 sancti Columbáni discípulus et ferme sexcentórum Monachórum Pater fuit; ac, 
 vitæ sanctitáte conspícuus, étiam miráculis 
 cláruit.
\switchcolumn
\selectlanguage{english}
In the monastery of Luxeuil, the death of Abbot St. Eustasius, a disciple of 
 St. Columban, who had under his guidance nearly six hundred monks. 
 Eminent in sanctity, he was also renowned for miracles.
\switchcolumn*
\selectlanguage{latin}
\end{paracol}


% ---- martyrology/mart03/mart0330.htm
\needspace{10\baselineskip}
\begin{paracol}{2}
\selectlanguage{latin}
\begin{center}{\color{gregoriocolor} Tértio Kaléndas Aprílis. 
 Luna\dots\ }\end{center}
\switchcolumn
\selectlanguage{english}
\begin{center}{\color{gregoriocolor} The 30\textsuperscript{th} Day of March. 
 The\dots\ Day of the Moon.}\end{center}
\end{paracol}

\noindent\begin{tabularx}{\linewidth}{*{19}{>{\centering\arraybackslash}X}}
 \textcolor{gregoriocolor}{a} & \textcolor{gregoriocolor}{b} & \textcolor{gregoriocolor}{c} & \textcolor{gregoriocolor}{d} & \textcolor{gregoriocolor}{e} & \textcolor{gregoriocolor}{f} & \textcolor{gregoriocolor}{g} & \textcolor{gregoriocolor}{h} & \textcolor{gregoriocolor}{i} & \textcolor{gregoriocolor}{k} & \textcolor{gregoriocolor}{l} & \textcolor{gregoriocolor}{m} & \textcolor{gregoriocolor}{n} & \textcolor{gregoriocolor}{p} & \textcolor{gregoriocolor}{q} & \textcolor{gregoriocolor}{r} & \textcolor{gregoriocolor}{s} & \textcolor{gregoriocolor}{t} & \textcolor{gregoriocolor}{u} \\
 1 & 2 & 3 & 4 & 5 & 6 & 7 & 8 & 9 & 10 & 11 & 12 & 13 & 14 & 15 & 16 & 17 & 18 & 19 \\
\end{tabularx}
\vspace{0.5\baselineskip}
\noindent\begin{tabularx}{\linewidth}{*{12}{>{\centering\arraybackslash}X}}
 \textcolor{gregoriocolor}{A} & \textcolor{gregoriocolor}{B} & \textcolor{gregoriocolor}{C} & \textcolor{gregoriocolor}{D} & \textcolor{gregoriocolor}{E} & F & \textcolor{gregoriocolor}{F} & \textcolor{gregoriocolor}{G} & \textcolor{gregoriocolor}{H} & \textcolor{gregoriocolor}{M} & \textcolor{gregoriocolor}{N} & \textcolor{gregoriocolor}{P} \\
 20 & 21 & 22 & 23 & 24 & 25 & 25 & 26 & 27 & 28 & 29 & 30 \\
\end{tabularx}

\begin{paracol}{2}
\selectlanguage{latin}
\lettrine[lines=2]{R}{omæ,} via Appia, pássio beáti Quiríni Tribúni, 
 patris sanctæ Balbínæ Vírginis, qui a beáto Alexándro Papa, quem habébat in 
 custódia, cum omni domo sua baptizátus est; atque, sub Hadriáno Imperatóre, 
 cum esset tráditus Aureliáno Júdici, et in fídei confessióne persísteret, 
 invíctus Christi miles, post linguæ abscissiónem, equúlei suspensiónem, 
 manuúmque ac pedum detruncatiónem, martyrii agónem gládio consummávit.
\switchcolumn
\selectlanguage{english}
\lettrine[lines=2]{A}{t} Rome, on the Appian Way, the martyrdom of the tribune blessed Quirinus, 
 who had been baptized with all his household by Pope St. Alexander when he 
 was imprisoned in their house. Under Emperor Adrian, he was delivered 
 to the judge Aurelian, and because he persevered in the confession of faith, 
 his tongue was torn out, he was stretched on the rack, his hands and feet 
 were cut off, and the sword completed his course of martyrdom.
\switchcolumn*
\selectlanguage{latin}
Thessalonícæ 
 natális sanctórum Mártyrum Domníni, Victóris et Sociórum.
\switchcolumn
\selectlanguage{english}
At Thessalonica, the birthday of the holy martyrs Domninus, Victor, and 
 their companions.
\switchcolumn*
\selectlanguage{latin}
Constantinópoli commemorátio sanctórum plurimórum Mártyrum cathólicæ 
 communiónis, quos, Constántii témpore, Macedónius hæresiárcha, inaudítis 
 tormentórum genéribus cruciátos, occídit; nam, inter cétera, fidélium 
 mulíerum úbera inter compréssa arcárum labra dissécuit, et candénti ferro 
 combússit.
\switchcolumn
\selectlanguage{english}
At Constantinople, in the time of Constantius, the commemoration of many 
 holy martyrs of the Catholic communion, whom the heresiarch Macedonius put 
 to death by unheard-of kinds of torments. Among other tortures, they 
 were burned with red-hot irons, and the breasts of Christian women were cut 
 away between the lids of coffers.
\switchcolumn*
\selectlanguage{latin}
In castro Silvanecténsi, in Gállia, deposítio sancti Réguli, Arelaténsis 
 Epíscopi.
\switchcolumn
\selectlanguage{english}
At Senlis in France, the death of St. Regulus, bishop of Arles.
\switchcolumn*
\selectlanguage{latin}
Auréliæ, in Gállia, sancti Pastóris Epíscopi.
\switchcolumn
\selectlanguage{english}
At Orleans in France, Bishop St. Pastor.
\switchcolumn*
\selectlanguage{latin}
Syracúsis, in Sicília, sancti Zósimi, 
 Epíscopi et Confessóris.
\switchcolumn
\selectlanguage{english}
At Syracuse, St. Zosimus, bishop and confessor.
\switchcolumn*
\selectlanguage{latin}
In monte Sina sancti Joánnis Clímaci Abbátis.
\switchcolumn
\selectlanguage{english}
On Mount Sinai, Abbot St. John Climacus.
\switchcolumn*
\selectlanguage{latin}
Aquilériæ, in Hispánia, sancti Petri Regaláti, 
 in urbe Vallisoletána orti, Sacerdótis ex Ordine Minórum et Confessóris, 
 reguláris disciplínæ in Hispániæ cœnóbiis restitutóris; quem Benedíctus 
 Décimus quartus, Póntifex Máximus, Sanctórum fastis adscrípsit.
\switchcolumn
\selectlanguage{english}
At Aquileria in Spain, the confessor St. Peter Regalado, priest of the 
 Order of Friars Minor. He was born in Valladolid, and restored the 
 regular discipline in the Spanish monasteries. Pope Benedict XIV 
 placed him on the roll of saints.
\switchcolumn*
\selectlanguage{latin}
Apud Aquínum sancti Clínii Confessóris.
\switchcolumn
\selectlanguage{english}
At Aquino, St. Clinius confessor.
\switchcolumn*
\selectlanguage{latin}
\end{paracol}


% ---- martyrology/mart03/mart0331.htm
\needspace{10\baselineskip}
\begin{paracol}{2}
\selectlanguage{latin}
\begin{center}{\color{gregoriocolor} Prídie Kaléndas Aprílis. 
 Luna\dots\ }\end{center}
\switchcolumn
\selectlanguage{english}
\begin{center}{\color{gregoriocolor} The 31\textsuperscript{st} Day of March. 
 The\dots\ Day of the Moon.}\end{center}
\end{paracol}

\noindent\begin{tabularx}{\linewidth}{*{19}{>{\centering\arraybackslash}X}}
 \textcolor{gregoriocolor}{a} & \textcolor{gregoriocolor}{b} & \textcolor{gregoriocolor}{c} & \textcolor{gregoriocolor}{d} & \textcolor{gregoriocolor}{e} & \textcolor{gregoriocolor}{f} & \textcolor{gregoriocolor}{g} & \textcolor{gregoriocolor}{h} & \textcolor{gregoriocolor}{i} & \textcolor{gregoriocolor}{k} & \textcolor{gregoriocolor}{l} & \textcolor{gregoriocolor}{m} & \textcolor{gregoriocolor}{n} & \textcolor{gregoriocolor}{p} & \textcolor{gregoriocolor}{q} & \textcolor{gregoriocolor}{r} & \textcolor{gregoriocolor}{s} & \textcolor{gregoriocolor}{t} & \textcolor{gregoriocolor}{u} \\
 2 & 3 & 4 & 5 & 6 & 7 & 8 & 9 & 10 & 11 & 12 & 13 & 14 & 15 & 16 & 17 & 18 & 19 & 20 \\
\end{tabularx}
\vspace{0.5\baselineskip}
\noindent\begin{tabularx}{\linewidth}{*{12}{>{\centering\arraybackslash}X}}
 \textcolor{gregoriocolor}{A} & \textcolor{gregoriocolor}{B} & \textcolor{gregoriocolor}{C} & \textcolor{gregoriocolor}{D} & \textcolor{gregoriocolor}{E} & F & \textcolor{gregoriocolor}{F} & \textcolor{gregoriocolor}{G} & \textcolor{gregoriocolor}{H} & \textcolor{gregoriocolor}{M} & \textcolor{gregoriocolor}{N} & \textcolor{gregoriocolor}{P} \\
 21 & 22 & 23 & 24 & 25 & 26 & 26 & 27 & 28 & 29 & 30 & 1 \\
\end{tabularx}

\begin{paracol}{2}
\selectlanguage{latin}
\lettrine[lines=2]{T}{hécuæ,} in Palæstína, sancti Amos Prophétæ, qui 
 ab Amasía Sacerdóte frequénter plagis afflíctus est, atque ab hujus fílio 
 Ozía vecte per témpora transfíxus; et póstea, semivívus in pátriam devéctus, 
 ibídem exspirávit, sepultúsque est cum pátribus suis.
\switchcolumn
\selectlanguage{english}
\lettrine[lines=2]{A}{t} Thecua in Palestine, the holy prophet Amos, whom the priest Amasias 
 frequently had scourged. Ozias, that priest's son, pierced his head at 
 the temples with an iron spike. Being carried half dead to his own 
 country, he died there, and was buried with his family.
\switchcolumn*
\selectlanguage{latin}
In Pérside sancti Bénjamin Diáconi, qui, cum Dei verbum non desísteret prædicáre, 
 ídeo, sub Isdegérde Rege, arundínibus acútis confíxus únguibus, et spinósa 
 sude per alvum transmíssa, martyrium consummávit.
\switchcolumn
\selectlanguage{english}
In Persia, during the reign of King Isdegerdes, the deacon St. Benjamin. 
 Because he would not stop preaching the word of God, he had a sharp reed 
 forced under his nails, a thorny stake driven through his body, and thus 
 completed his martyrdom.
\switchcolumn*
\selectlanguage{latin}
In Africa sanctórum Mártyrum Theodúli, Anésii, Felícis, Cornéliæ 
 et Sociórum.
\switchcolumn
\selectlanguage{english}
In Africa, the holy martyrs Theodulus, Anesius, Felix, Cornelia, and their 
 companions.
\switchcolumn*
\selectlanguage{latin}
Romæ sanctæ Balbínæ Vírginis, fíliæ beáti 
 Quiríni Mártyris, quæ, a sancto Alexándro Papa baptizáta, in sancta 
 virginitáte Christum sibi sponsum elégit; et, post devíctum hujus sæculi 
 cursum, sepúlta est via Appia, juxta patrem suum.
\switchcolumn
\selectlanguage{english}
At Rome, the virgin St. Balbina, daughter of the blessed martyr Quirinus. 
 She was baptized by Pope Alexander, and she chose Christ as her spouse in 
 her virginity. After overcoming the world, she was buried at her 
 father's side on the Appian Way.
\switchcolumn*
\selectlanguage{latin}
\end{paracol}
